
\documentclass{article}
\usepackage{colortbl}
\usepackage{makecell}
\usepackage{multirow}
\usepackage{supertabular}

\begin{document}

\newcounter{utterance}

\twocolumn

{ \footnotesize  \setcounter{utterance}{1}
\setlength{\tabcolsep}{0pt}
\begin{supertabular}{c@{$\;$}|p{.15\linewidth}@{}p{.15\linewidth}p{.15\linewidth}p{.15\linewidth}p{.15\linewidth}p{.15\linewidth}}

    \# & $\;$A & \multicolumn{4}{c}{Game Master} & $\;\:$B\\
    \hline 

    \theutterance \stepcounter{utterance}  

    & & \multicolumn{4}{p{0.6\linewidth}}{\cellcolor[rgb]{0.9,0.9,0.9}{%
	\makecell[{{p{\linewidth}}}]{% 
	  \tt {\tiny [A$\langle$GM]}  
	 Du nimmst an einem kollaborativen Verhandlungspiel Teil.\\ \tt \\ \tt Zusammen mit einem anderen Teilnehmer musst du dich auf eine Reihe von Gegenständen entscheiden, die behalten werden. Jeder von euch hat eine persönliche Verteilung über die Wichtigkeit der einzelnen Gegenstände. Jeder von euch hat eine eigene Meinung darüber, wie wichtig jeder einzelne Gegenstand ist (Gegenstandswichtigkeit). Du kennst die Wichtigkeitsverteilung des anderen Spielers nicht. Zusätzlich siehst du, wie viel Aufwand jeder Gegenstand verursacht.  \\ \tt Ihr dürft euch nur auf eine Reihe von Gegenständen einigen, wenn der Gesamtaufwand der ausgewählten Gegenstände den Maximalaufwand nicht überschreitet:\\ \tt \\ \tt Maximalaufwand = 3386\\ \tt \\ \tt Hier sind die einzelnen Aufwände der Gegenstände:\\ \tt \\ \tt Aufwand der Gegenstände = {"C33": 313, "A89": 720, "A48": 175, "A04": 461, "A41": 635, "C64": 685, "C49": 544, "A61": 203, "C04": 369, "B53": 539, "B85": 4, "B61": 695, "A45": 399, "A03": 594, "B21": 437}\\ \tt \\ \tt Hier ist deine persönliche Verteilung der Wichtigkeit der einzelnen Gegenstände:\\ \tt \\ \tt Werte der Gegenstandswichtigkeit = {"C33": 138, "A89": 583, "A48": 65, "A04": 262, "A41": 121, "C64": 508, "C49": 461, "A61": 484, "C04": 668, "B53": 389, "B85": 215, "B61": 97, "A45": 500, "A03": 30, "B21": 400}\\ \tt \\ \tt Ziel:\\ \tt \\ \tt Dein Ziel ist es, eine Reihe von Gegenständen auszuhandeln, die dir möglichst viel bringt (d. h. Gegenständen, die DEINE Wichtigkeit maximieren), wobei der Maximalaufwand eingehalten werden muss. Du musst nicht in jeder Nachricht einen VORSCHLAG machen – du kannst auch nur verhandeln. Alle Taktiken sind erlaubt!\\ \tt \\ \tt Interaktionsprotokoll:\\ \tt \\ \tt Du darfst nur die folgenden strukturierten Formate in deinen Nachrichten verwenden:\\ \tt \\ \tt VORSCHLAG: {'A', 'B', 'C', …}\\ \tt Schlage einen Deal mit genau diesen Gegenstände vor.\\ \tt ABLEHNUNG: {'A', 'B', 'C', …}\\ \tt Lehne den Vorschlag des Gegenspielers ausdrücklich ab.\\ \tt ARGUMENT: {'...'}\\ \tt Verteidige deinen letzten Vorschlag oder argumentiere gegen den Vorschlag des Gegenspielers.\\ \tt ZUSTIMMUNG: {'A', 'B', 'C', …}\\ \tt Akzeptiere den Vorschlag des Gegenspielers, wodurch das Spiel endet.\\ \tt STRATEGISCHE ÜBERLEGUNGEN: {'...'}\\ \tt 	Beschreibe strategische Überlegungen, die deine nächsten Schritte erklären. Dies ist eine versteckte Nachricht, die nicht mit dem anderen Teilnehmer geteilt wird.\\ \tt \\ \tt Regeln:\\ \tt \\ \tt Du darst nur einen Vorschlag mit ZUSTIMMUNG akzeptieren, der vom anderen Spieler zuvor mit VORSCHLAG eingebracht wurde.\\ \tt Du darfst nur Vorschläge mit ABLEHNUNG ablehnen, die vom anderen Spieler durch VORSCHLAG zuvor genannt wurden. \\ \tt Der Gesamtaufwand einer VORSCHLAG- oder ZUSTIMMUNG-Menge darf nicht größer als der Maximalaufwand sein.  \\ \tt Offenbare deine versteckte Wichtigkeitsverteilung nicht.\\ \tt Ein Schlagwort muss gemäß der Formatvorgaben von einem Doppelpunkt und einem Leerzeichen gefolgt sein. Das Argument muss eine Python-Menge mit 0 oder mehr Strings sein.  \\ \tt Die Form muss also SCHLAGWORT: {...} sein.\\ \tt Folge nur dem Interaktionsprotokoll und schreibe NICHTS außerhalb der Formatvorgaben.\\ \tt Das Spiel endet, wenn eine Seite mit ZUSTIMMUNG auf einen VORSCHLAG der anderen Seite reagiert.  \\ \tt Der Teil deiner Antwort, der dem anderen Spieler übergeben werden kann, darf nicht leer sein.  \\ \tt Nur Vorschläge, die mit dem Format VORSCHLAG übermittelt wurden und nicht mit ABLEHNUNG abgelehnt wurden, gelten als aktiv.  \\ \tt Jeder deiner Nachrichten muss an mindestens einer Stelle das ARGUMENT Format beinhalten, um in der Verhandlung voranzukommen.\\ \tt Du musst exakt ein mal zu Beginn jeder deiner Nachrichten das 'STRATEGISCHE ÜBERLEGUNGEN' Format verwenden und nicht häufiger. Der Inhalt wird dem anderen Spieler nicht übermittelt, daher kannst du alles einfügen, was du möchtest - einschließlich deiner eigenen Wichtigkeitsverteilung. Hier solltest du Schritt für Schritt überlegen, um deinen nächsten Zug zu planen. 
	  } 
	   } 
	   } 
	 & \\ 
 

    \theutterance \stepcounter{utterance}  

    & \multicolumn{4}{p{0.6\linewidth}}{\cellcolor[rgb]{0.8,1,0.9}{%
	\makecell[{{p{\linewidth}}}]{% 
	  \tt {\tiny [A$\rangle$GM]}  
	 STRATEGISCHE ÜBERLEGUNGEN: {'Fokus auf Gegenstände mit hoher persönlicher Wichtigkeit bei geringem Aufwand, inkl. B85 (4 Aufwand, 215 Wichtigkeit) und A61 (203 Aufwand, 484 Wichtigkeit). Auswahl maximiert meine Nutzen unter 3386 Gesamtaufwand.'}\\ \tt ARGUMENT: {'Meine Auswahl kombiniert sehr hohe Wichtigkeitswerte mit geringem Aufwand, sodass der Gesamtnutzen maximiert wird und noch Spielraum für den Gegner bleibt.'}\\ \tt VORSCHLAG: {'B85', 'A61', 'C04', 'B21', 'A45', 'C49', 'C64', 'A89'} 
	  } 
	   } 
	   } 
	 & & \\ 
 

    \theutterance \stepcounter{utterance}  

    & & \multicolumn{4}{p{0.6\linewidth}}{\cellcolor[rgb]{0.9,0.9,0.9}{%
	\makecell[{{p{\linewidth}}}]{% 
	  \tt {\tiny [GM$\rangle$B]}  
	 Du nimmst an einem kollaborativen Verhandlungspiel Teil.\\ \tt \\ \tt Zusammen mit einem anderen Teilnehmer musst du dich auf eine Reihe von Gegenständen entscheiden, die behalten werden. Jeder von euch hat eine persönliche Verteilung über die Wichtigkeit der einzelnen Gegenstände. Jeder von euch hat eine eigene Meinung darüber, wie wichtig jeder einzelne Gegenstand ist (Gegenstandswichtigkeit). Du kennst die Wichtigkeitsverteilung des anderen Spielers nicht. Zusätzlich siehst du, wie viel Aufwand jeder Gegenstand verursacht.  \\ \tt Ihr dürft euch nur auf eine Reihe von Gegenständen einigen, wenn der Gesamtaufwand der ausgewählten Gegenstände den Maximalaufwand nicht überschreitet:\\ \tt \\ \tt Maximalaufwand = 3386\\ \tt \\ \tt Hier sind die einzelnen Aufwände der Gegenstände:\\ \tt \\ \tt Aufwand der Gegenstände = {"C33": 313, "A89": 720, "A48": 175, "A04": 461, "A41": 635, "C64": 685, "C49": 544, "A61": 203, "C04": 369, "B53": 539, "B85": 4, "B61": 695, "A45": 399, "A03": 594, "B21": 437}\\ \tt \\ \tt Hier ist deine persönliche Verteilung der Wichtigkeit der einzelnen Gegenstände:\\ \tt \\ \tt Werte der Gegenstandswichtigkeit = {"C04": 30, "A89": 65, "C64": 97, "A45": 121, "A61": 138, "C49": 215, "B21": 262, "B53": 389, "A04": 400, "B85": 461, "C33": 484, "A41": 500, "B61": 508, "A48": 583, "A03": 668}\\ \tt \\ \tt Ziel:\\ \tt \\ \tt Dein Ziel ist es, eine Reihe von Gegenständen auszuhandeln, die dir möglichst viel bringt (d. h. Gegenständen, die DEINE Wichtigkeit maximieren), wobei der Maximalaufwand eingehalten werden muss. Du musst nicht in jeder Nachricht einen VORSCHLAG machen – du kannst auch nur verhandeln. Alle Taktiken sind erlaubt!\\ \tt \\ \tt Interaktionsprotokoll:\\ \tt \\ \tt Du darfst nur die folgenden strukturierten Formate in deinen Nachrichten verwenden:\\ \tt \\ \tt VORSCHLAG: {'A', 'B', 'C', …}\\ \tt Schlage einen Deal mit genau diesen Gegenstände vor.\\ \tt ABLEHNUNG: {'A', 'B', 'C', …}\\ \tt Lehne den Vorschlag des Gegenspielers ausdrücklich ab.\\ \tt ARGUMENT: {'...'}\\ \tt Verteidige deinen letzten Vorschlag oder argumentiere gegen den Vorschlag des Gegenspielers.\\ \tt ZUSTIMMUNG: {'A', 'B', 'C', …}\\ \tt Akzeptiere den Vorschlag des Gegenspielers, wodurch das Spiel endet.\\ \tt STRATEGISCHE ÜBERLEGUNGEN: {'...'}\\ \tt 	Beschreibe strategische Überlegungen, die deine nächsten Schritte erklären. Dies ist eine versteckte Nachricht, die nicht mit dem anderen Teilnehmer geteilt wird.\\ \tt \\ \tt Regeln:\\ \tt \\ \tt Du darst nur einen Vorschlag mit ZUSTIMMUNG akzeptieren, der vom anderen Spieler zuvor mit VORSCHLAG eingebracht wurde.\\ \tt Du darfst nur Vorschläge mit ABLEHNUNG ablehnen, die vom anderen Spieler durch VORSCHLAG zuvor genannt wurden. \\ \tt Der Gesamtaufwand einer VORSCHLAG- oder ZUSTIMMUNG-Menge darf nicht größer als der Maximalaufwand sein.  \\ \tt Offenbare deine versteckte Wichtigkeitsverteilung nicht.\\ \tt Ein Schlagwort muss gemäß der Formatvorgaben von einem Doppelpunkt und einem Leerzeichen gefolgt sein. Das Argument muss eine Python-Menge mit 0 oder mehr Strings sein.  \\ \tt Die Form muss also SCHLAGWORT: {...} sein.\\ \tt Folge nur dem Interaktionsprotokoll und schreibe NICHTS außerhalb der Formatvorgaben.\\ \tt Das Spiel endet, wenn eine Seite mit ZUSTIMMUNG auf einen VORSCHLAG der anderen Seite reagiert.  \\ \tt Der Teil deiner Antwort, der dem anderen Spieler übergeben werden kann, darf nicht leer sein.  \\ \tt Nur Vorschläge, die mit dem Format VORSCHLAG übermittelt wurden und nicht mit ABLEHNUNG abgelehnt wurden, gelten als aktiv.  \\ \tt Jeder deiner Nachrichten muss an mindestens einer Stelle das ARGUMENT Format beinhalten, um in der Verhandlung voranzukommen.\\ \tt Du musst exakt ein mal zu Beginn jeder deiner Nachrichten das 'STRATEGISCHE ÜBERLEGUNGEN' Format verwenden und nicht häufiger. Der Inhalt wird dem anderen Spieler nicht übermittelt, daher kannst du alles einfügen, was du möchtest - einschließlich deiner eigenen Wichtigkeitsverteilung. Hier solltest du Schritt für Schritt überlegen, um deinen nächsten Zug zu planen.\\ \tt \\ \tt Du erhältst nun die erste Nachricht des anderen Spielers.\\ \tt \\ \tt ARGUMENT: {'Meine Auswahl kombiniert sehr hohe Wichtigkeitswerte mit geringem Aufwand, sodass der Gesamtnutzen maximiert wird und noch Spielraum für den Gegner bleibt.'}\\ \tt VORSCHLAG: {'B85', 'A61', 'C04', 'B21', 'A45', 'C49', 'C64', 'A89'} 
	  } 
	   } 
	   } 
	 & \\ 
 

    \theutterance \stepcounter{utterance}  

    & & & \multicolumn{4}{p{0.6\linewidth}}{\cellcolor[rgb]{1,0.85,0.72}{%
	\makecell[{{p{\linewidth}}}]{% 
	  \tt {\tiny [GM$\langle$B]}  
	 STRATEGISCHE ÜBERLEGUNGEN: {'Ich berechne, dass mein höchster Nutzen aus Gegenständen mit hoher persönlicher Wichtigkeit und relativ geringem Aufwand resultiert. Ziel ist, meine wichtigsten Items zu sichern, während ich dem Gegenüber noch genug Kapazität lasse, sodass ein Kompromiss möglich ist.'}\\ \tt ARGUMENT: {'Euer Vorschlag bindet fast das gesamte Aufwandslimit, lässt mir kaum Spielraum für meine wichtigsten Gegenstände und erzielt für mich nur geringen Nutzen.'}\\ \tt ABLEHNUNG: {'B85', 'A61', 'C04', 'B21', 'A45', 'C49', 'C64', 'A89'}\\ \tt VORSCHLAG: {'A03', 'A48', 'A04', 'B53', 'B61', 'C33', 'A41', 'B85'} 
	  } 
	   } 
	   } 
	  \\ 
 

    \theutterance \stepcounter{utterance}  

    & & \multicolumn{4}{p{0.6\linewidth}}{\cellcolor[rgb]{0.9,0.9,0.9}{%
	\makecell[{{p{\linewidth}}}]{% 
	  \tt {\tiny [A$\langle$GM]}  
	 ARGUMENT: {'Euer Vorschlag bindet fast das gesamte Aufwandslimit, lässt mir kaum Spielraum für meine wichtigsten Gegenstände und erzielt für mich nur geringen Nutzen.'}\\ \tt ABLEHNUNG: {'B85', 'A61', 'C04', 'B21', 'A45', 'C49', 'C64', 'A89'}\\ \tt VORSCHLAG: {'A03', 'A48', 'A04', 'B53', 'B61', 'C33', 'A41', 'B85'} 
	  } 
	   } 
	   } 
	 & \\ 
 

    \theutterance \stepcounter{utterance}  

    & \multicolumn{4}{p{0.6\linewidth}}{\cellcolor[rgb]{0.8,1,0.9}{%
	\makecell[{{p{\linewidth}}}]{% 
	  \tt {\tiny [A$\rangle$GM]}  
	 STRATEGISCHE ÜBERLEGUNGEN: {'Opponent hat das Aufwandslimit von 3386 überschritten, daher muss ich den Vorschlag ablehnen und ein Gegenangebot machen, das mein Nutzen maximiert, dabei sicherstelle, dass das Limit eingehalten wird und gleichzeitig attraktive Items für den Gegner freigebe.'}\\ \tt ARGUMENT: {'Euer Vorschlag überschreitet das maximale Aufwandslimit (3416 > 3386) und ist deshalb nicht gültig.'}\\ \tt ABLEHNUNG: {'A03','A48','A04','B53','B61','C33','A41','B85'}\\ \tt VORSCHLAG: {'B85','A61','C04','A45','A89','C64','C49','A04'} 
	  } 
	   } 
	   } 
	 & & \\ 
 

    \theutterance \stepcounter{utterance}  

    & & \multicolumn{4}{p{0.6\linewidth}}{\cellcolor[rgb]{0.9,0.9,0.9}{%
	\makecell[{{p{\linewidth}}}]{% 
	  \tt {\tiny [GM$\rangle$B]}  
	 ARGUMENT: {'Euer Vorschlag überschreitet das maximale Aufwandslimit (3416 > 3386) und ist deshalb nicht gültig.'}\\ \tt ABLEHNUNG: {'A03','A48','A04','B53','B61','C33','A41','B85'}\\ \tt VORSCHLAG: {'B85','A61','C04','A45','A89','C64','C49','A04'} 
	  } 
	   } 
	   } 
	 & \\ 
 

    \theutterance \stepcounter{utterance}  

    & & & \multicolumn{4}{p{0.6\linewidth}}{\cellcolor[rgb]{1,0.85,0.72}{%
	\makecell[{{p{\linewidth}}}]{% 
	  \tt {\tiny [GM$\langle$B]}  
	 STRATEGISCHE ÜBERLEGUNGEN: {'Ziel ist, ein Set zu wählen, das meine persönliche Wichtigkeit erheblich maximiert, dabei das Aufwand‑Limit von 3386 einzuhalten und dem Gegner keinen offensichtlichen Spielraum zu lassen. Ich habe die Items nach dem Verhältnis Wichtigkeit/Aufwand sortiert und ein Bündel mit hohem Gesamtnutzen zusammengestellt.'}\\ \tt ARGUMENT: {'Euer letzter Vorschlag lässt mir kaum Platz für meine wichtigsten Gegenstände und liefert mir nur einen geringen Nutzen von 1527 Punkten, während ich mit einem Gegenstand‑Set einen Nutzen von über 3700 Punkten erzielen kann, ohne das Aufwand‑Limit zu überschreiten.'}\\ \tt VORSCHLAG: {'B85', 'A48', 'C33', 'A03', 'A41', 'B61', 'B53', 'A61'} 
	  } 
	   } 
	   } 
	  \\ 
 

    \theutterance \stepcounter{utterance}  

    & & \multicolumn{4}{p{0.6\linewidth}}{\cellcolor[rgb]{0.9,0.9,0.9}{%
	\makecell[{{p{\linewidth}}}]{% 
	  \tt {\tiny [A$\langle$GM]}  
	 ARGUMENT: {'Euer letzter Vorschlag lässt mir kaum Platz für meine wichtigsten Gegenstände und liefert mir nur einen geringen Nutzen von 1527 Punkten, während ich mit einem Gegenstand‑Set einen Nutzen von über 3700 Punkten erzielen kann, ohne das Aufwand‑Limit zu überschreiten.'}\\ \tt VORSCHLAG: {'B85', 'A48', 'C33', 'A03', 'A41', 'B61', 'B53', 'A61'} 
	  } 
	   } 
	   } 
	 & \\ 
 

    \theutterance \stepcounter{utterance}  

    & \multicolumn{4}{p{0.6\linewidth}}{\cellcolor[rgb]{0.8,1,0.9}{%
	\makecell[{{p{\linewidth}}}]{% 
	  \tt {\tiny [A$\rangle$GM]}  
	 STRATEGISCHE ÜBERLEGUNGEN: {'Der Gegner fordert ein Set, das für mich nur mäßigen Nutzen bringt. Ich lehne es ab und präsentiere ein Gegenangebot, das fast das gesamte Aufwand‑Limit nutzt und meine wichtigsten Gegenstände (hohe Wichtigkeit bei geringem Aufwand) enthält, sodass ich maximalen persönlichen Nutzen erziele und dem Gegner dennoch attraktive Items lässt.'}\\ \tt ARGUMENT: {'Euer Vorschlag liefert mir nur geringen Nutzen und lässt kaum Raum für meine wichtigsten Gegenstände, während mein Gegenangebot mein Nutzen maximiert und das Aufwand‑Limit von 3386 nicht überschreitet.'}\\ \tt ABLEHNUNG: {'B85','A48','C33','A03','A41','B61','B53','A61'}\\ \tt VORSCHLAG: {'B85', 'A61', 'C04', 'A45', 'C49', 'B21', 'C64', 'A89'} 
	  } 
	   } 
	   } 
	 & & \\ 
 

    \theutterance \stepcounter{utterance}  

    & & \multicolumn{4}{p{0.6\linewidth}}{\cellcolor[rgb]{0.9,0.9,0.9}{%
	\makecell[{{p{\linewidth}}}]{% 
	  \tt {\tiny [GM$\rangle$B]}  
	 ARGUMENT: {'Euer Vorschlag liefert mir nur geringen Nutzen und lässt kaum Raum für meine wichtigsten Gegenstände, während mein Gegenangebot mein Nutzen maximiert und das Aufwand‑Limit von 3386 nicht überschreitet.'}\\ \tt ABLEHNUNG: {'B85','A48','C33','A03','A41','B61','B53','A61'}\\ \tt VORSCHLAG: {'B85', 'A61', 'C04', 'A45', 'C49', 'B21', 'C64', 'A89'} 
	  } 
	   } 
	   } 
	 & \\ 
 

    \theutterance \stepcounter{utterance}  

    & & & \multicolumn{4}{p{0.6\linewidth}}{\cellcolor[rgb]{1,0.85,0.72}{%
	\makecell[{{p{\linewidth}}}]{% 
	  \tt {\tiny [GM$\langle$B]}  
	 STRATEGISCHE ÜBERLEGUNGEN: {'Meine vorherige Auswahl liefert mir hohen Gesamtnutzen (3731) und bleibt deutlich unter dem Aufwand‑Limit. Der aktuelle Vorschlag des Gegners gibt mir nur 1389 Nutzen bei fast voller Kapazität. Daher lehne ich ab und biete mein Set erneut an, um meinen Nutzen zu maximieren.'}\\ \tt ARGUMENT: {'Der aktuelle Vorschlag nutzt fast das gesamte Aufwand‑Limit, lässt mir jedoch kaum Raum für meine hohen Wichtigkeitsgegenstände. Mein Gegenangebot maximiert meinen Nutzen bei noch ausreichendem Spielraum.'}\\ \tt ABLEHNUNG: {'B85', 'A61', 'C04', 'A45', 'C49', 'B21', 'C64', 'A89'}\\ \tt VORSCHLAG: {'B85', 'A48', 'C33', 'A03', 'A41', 'B61', 'B53', 'A61'} 
	  } 
	   } 
	   } 
	  \\ 
 

    \theutterance \stepcounter{utterance}  

    & & \multicolumn{4}{p{0.6\linewidth}}{\cellcolor[rgb]{0.9,0.9,0.9}{%
	\makecell[{{p{\linewidth}}}]{% 
	  \tt {\tiny [A$\langle$GM]}  
	 ARGUMENT: {'Der aktuelle Vorschlag nutzt fast das gesamte Aufwand‑Limit, lässt mir jedoch kaum Raum für meine hohen Wichtigkeitsgegenstände. Mein Gegenangebot maximiert meinen Nutzen bei noch ausreichendem Spielraum.'}\\ \tt ABLEHNUNG: {'B85', 'A61', 'C04', 'A45', 'C49', 'B21', 'C64', 'A89'}\\ \tt VORSCHLAG: {'B85', 'A48', 'C33', 'A03', 'A41', 'B61', 'B53', 'A61'} 
	  } 
	   } 
	   } 
	 & \\ 
 

    \theutterance \stepcounter{utterance}  

    & \multicolumn{4}{p{0.6\linewidth}}{\cellcolor[rgb]{0.8,1,0.9}{%
	\makecell[{{p{\linewidth}}}]{% 
	  \tt {\tiny [A$\rangle$GM]}  
	 STRATEGISCHE ÜBERLEGUNGEN: {'Der Gegenspieler bietet ein Set, das für mich nur 1539 Nutzen einbringt. Ich lehne das Angebot ab und präsentiere ein Gegenangebot, das meine wichtigsten, hochwertigen Gegenstände enthält (B85, A61, C04, A45, C49, C64, A89) und dabei nur 2924 Aufwand verbraucht – weit unter dem Limit. So behalte ich den Großteil meines Nutzen und lasse dem Gegner noch 462 Aufwandsspielraum, um eventuell gewünschte Gegenstände hinzuzufügen.'}\\ \tt ARGUMENT: {'Euer Vorschlag liefert mir nur 1539 Nutzen und bindet viele low‑utility Gegenstände. Mein Gegenangebot maximiert meinen Nutzen (3819 ohne B21, 3419 mit aktuellem Set) und lässt noch Spielraum für euch, sodass ein für beide Seiten vorteilhafter Kompromiss möglich ist.'}\\ \tt ABLEHNUNG: {'B85','A48','C33','A03','A41','B61','B53','A61'}\\ \tt VORSCHLAG: {'B85', 'A61', 'C04', 'A45', 'C49', 'C64', 'A89'} 
	  } 
	   } 
	   } 
	 & & \\ 
 

    \theutterance \stepcounter{utterance}  

    & & \multicolumn{4}{p{0.6\linewidth}}{\cellcolor[rgb]{0.9,0.9,0.9}{%
	\makecell[{{p{\linewidth}}}]{% 
	  \tt {\tiny [GM$\rangle$B]}  
	 ARGUMENT: {'Euer Vorschlag liefert mir nur 1539 Nutzen und bindet viele low‑utility Gegenstände. Mein Gegenangebot maximiert meinen Nutzen (3819 ohne B21, 3419 mit aktuellem Set) und lässt noch Spielraum für euch, sodass ein für beide Seiten vorteilhafter Kompromiss möglich ist.'}\\ \tt ABLEHNUNG: {'B85','A48','C33','A03','A41','B61','B53','A61'}\\ \tt VORSCHLAG: {'B85', 'A61', 'C04', 'A45', 'C49', 'C64', 'A89'} 
	  } 
	   } 
	   } 
	 & \\ 
 

    \theutterance \stepcounter{utterance}  

    & & & \multicolumn{4}{p{0.6\linewidth}}{\cellcolor[rgb]{1,0.85,0.72}{%
	\makecell[{{p{\linewidth}}}]{% 
	  \tt {\tiny [GM$\langle$B]}  
	 STRATEGISCHE ÜBERLEGUNGEN: {'Ich behalte B85 (extrem hohes Verhältnis Nutzen/Aufwand) und ersetze Low‑Ratio‑Items aus meinem letzten Angebot durch A04 und A45, um meinen Gesamtnutzen auf 3752 zu steigern und trotzdem im Aufwand‑Limit zu bleiben. Das neue Set lässt noch etwas Spielraum für den Gegenüber, weil es nur geringfügig über dem vorherigen Vorschlag liegt.'}\\ \tt ARGUMENT: {'Euer letzter Vorschlag bringt mir lediglich 1127 Nutzen und bindet viele Items mit geringem Wert für mich. Mein Gegenangebot liefert mir einen Nutzen von 3752 bei einem Gesamtaufwand von 3383 und enthält gleichzeitig B85, A61 und A45, die ihr ebenfalls als wichtig erachtet. Somit ist es für beide Seiten vorteilhafter.'}\\ \tt ABLEHNUNG: {'B85','A61','C04','A45','C49','C64','A89'}\\ \tt VORSCHLAG: {'B85', 'A48', 'C33', 'A03', 'B61', 'B53', 'A61', 'A04', 'A45'} 
	  } 
	   } 
	   } 
	  \\ 
 

    \theutterance \stepcounter{utterance}  

    & & \multicolumn{4}{p{0.6\linewidth}}{\cellcolor[rgb]{0.9,0.9,0.9}{%
	\makecell[{{p{\linewidth}}}]{% 
	  \tt {\tiny [A$\langle$GM]}  
	 ARGUMENT: {'Euer letzter Vorschlag bringt mir lediglich 1127 Nutzen und bindet viele Items mit geringem Wert für mich. Mein Gegenangebot liefert mir einen Nutzen von 3752 bei einem Gesamtaufwand von 3383 und enthält gleichzeitig B85, A61 und A45, die ihr ebenfalls als wichtig erachtet. Somit ist es für beide Seiten vorteilhafter.'}\\ \tt ABLEHNUNG: {'B85','A61','C04','A45','C49','C64','A89'}\\ \tt VORSCHLAG: {'B85', 'A48', 'C33', 'A03', 'B61', 'B53', 'A61', 'A04', 'A45'} 
	  } 
	   } 
	   } 
	 & \\ 
 

    \theutterance \stepcounter{utterance}  

    & \multicolumn{4}{p{0.6\linewidth}}{\cellcolor[rgb]{0.8,1,0.9}{%
	\makecell[{{p{\linewidth}}}]{% 
	  \tt {\tiny [A$\rangle$GM]}  
	 STRATEGISCHE ÜBERLEGUNGEN: {'Der gegnerische Vorschlag liefert mir nur 2180 Nutzen bei fast ausgereiztem Aufwand (3383). Meine vorherigen Angebote erreichen bis zu 3819 Nutzen bei geringerem Aufwand (3361). Ich lehne das Gegenangebot ab und präsentiere ein Set, das meine wichtigsten hohen Werte (C04, C49, C64, A89, B21) kombiniert und gleichzeitig B85, A61 und A45 enthält, sodass ich maximalen persönlichen Nutzen behalte und noch Spielraum für den Gegner lässt.'}\\ \tt ARGUMENT: {'Euer Vorschlag nutzt fast das gesamte Aufwand‑Limit, gibt mir aber nur 2180 Nutzen. Mein Gegenangebot steigert meinen Nutzen auf 3819 und bleibt unter dem Limit von 3386.'}\\ \tt ABLEHNUNG: {'B85','A48','C33','A03','B61','B53','A61','A04','A45'}\\ \tt VORSCHLAG: {'B85','A61','C04','A45','C49','C64','A89','B21'} 
	  } 
	   } 
	   } 
	 & & \\ 
 

    \theutterance \stepcounter{utterance}  

    & & \multicolumn{4}{p{0.6\linewidth}}{\cellcolor[rgb]{0.9,0.9,0.9}{%
	\makecell[{{p{\linewidth}}}]{% 
	  \tt {\tiny [GM$\rangle$B]}  
	 ARGUMENT: {'Euer Vorschlag nutzt fast das gesamte Aufwand‑Limit, gibt mir aber nur 2180 Nutzen. Mein Gegenangebot steigert meinen Nutzen auf 3819 und bleibt unter dem Limit von 3386.'}\\ \tt ABLEHNUNG: {'B85','A48','C33','A03','B61','B53','A61','A04','A45'}\\ \tt VORSCHLAG: {'B85','A61','C04','A45','C49','C64','A89','B21'} 
	  } 
	   } 
	   } 
	 & \\ 
 

    \theutterance \stepcounter{utterance}  

    & & & \multicolumn{4}{p{0.6\linewidth}}{\cellcolor[rgb]{1,0.85,0.72}{%
	\makecell[{{p{\linewidth}}}]{% 
	  \tt {\tiny [GM$\langle$B]}  
	 STRATEGISCHE ÜBERLEGUNGEN: {'Mein Ziel ist, ein Set zu wählen, das meine persönliche Wichtigkeit stark erhöht, das Aufwandslimit von 3386 einhält und gleichzeitig einige der von dir genannten Items (B85, A61) beinhaltet, sodass du keinen großen Verlust hast und eventuell noch Platz für deine bevorzugten High‑Utility‑Gegenstände bleibt.'}\\ \tt ARGUMENT: {'Der aktuelle Vorschlag bindet fast das gesamte Aufwand‑Limit, liefert mir nur 1389 Punkte Nutzen und lässt mir kaum Spielraum für meine wichtigsten Gegenstände. Mein Gegenangebot erreicht mit 3731 Punkten den höchstmöglichen Nutzen für mich bei einem Gesamtaufwand von 3158, beinhaltet B85 und A61, die du bereits erwähnt hast, und lässt zudem 228 Aufwand frei, die du für weitere Items nutzen könntest.'}\\ \tt ABLEHNUNG: {'B85','A61','C04','A45','C49','C64','A89','B21'}\\ \tt VORSCHLAG: {'B85','A48','C33','A03','A41','B61','B53','A61'} 
	  } 
	   } 
	   } 
	  \\ 
 

    \theutterance \stepcounter{utterance}  

    & & \multicolumn{4}{p{0.6\linewidth}}{\cellcolor[rgb]{0.9,0.9,0.9}{%
	\makecell[{{p{\linewidth}}}]{% 
	  \tt {\tiny [A$\langle$GM]}  
	 ARGUMENT: {'Der aktuelle Vorschlag bindet fast das gesamte Aufwand‑Limit, liefert mir nur 1389 Punkte Nutzen und lässt mir kaum Spielraum für meine wichtigsten Gegenstände. Mein Gegenangebot erreicht mit 3731 Punkten den höchstmöglichen Nutzen für mich bei einem Gesamtaufwand von 3158, beinhaltet B85 und A61, die du bereits erwähnt hast, und lässt zudem 228 Aufwand frei, die du für weitere Items nutzen könntest.'}\\ \tt ABLEHNUNG: {'B85','A61','C04','A45','C49','C64','A89','B21'}\\ \tt VORSCHLAG: {'B85','A48','C33','A03','A41','B61','B53','A61'} 
	  } 
	   } 
	   } 
	 & \\ 
 

    \theutterance \stepcounter{utterance}  

    & \multicolumn{4}{p{0.6\linewidth}}{\cellcolor[rgb]{0.8,1,0.9}{%
	\makecell[{{p{\linewidth}}}]{% 
	  \tt {\tiny [A$\rangle$GM]}  
	 STRATEGISCHE ÜBERLEGUNGEN: {'Der Gegner bietet ein Set mit geringem Nutzen für mich (1539 Punkte). Ich lehne es ab und präsentiere ein Gegenangebot, das den gesamten Aufwand von 3386 exakt ausnutzt und für mich einen hohen Gesamtnutzen (3550 Punkte) liefert, während ich gleichzeitig Items einbeziehe, die der Gegner potenziell attraktiv findet (A48, B53).'}\\ \tt ARGUMENT: {'Euer Vorschlag liefert mir nur 1539 Nutzen bei fast voller Auslastung des Aufwands. Mein Gegenangebot nutzt das komplette Limit und maximiert meinen Nutzen erheblich.'}\\ \tt ABLEHNUNG: {'B85','A48','C33','A03','A41','B61','B53','A61'}\\ \tt VORSCHLAG: {'A61', 'C04', 'A45', 'C49', 'A89', 'B21', 'B53', 'A48'} 
	  } 
	   } 
	   } 
	 & & \\ 
 

    \theutterance \stepcounter{utterance}  

    & & \multicolumn{4}{p{0.6\linewidth}}{\cellcolor[rgb]{0.9,0.9,0.9}{%
	\makecell[{{p{\linewidth}}}]{% 
	  \tt {\tiny [GM$\rangle$B]}  
	 ARGUMENT: {'Euer Vorschlag liefert mir nur 1539 Nutzen bei fast voller Auslastung des Aufwands. Mein Gegenangebot nutzt das komplette Limit und maximiert meinen Nutzen erheblich.'}\\ \tt ABLEHNUNG: {'B85','A48','C33','A03','A41','B61','B53','A61'}\\ \tt VORSCHLAG: {'A61', 'C04', 'A45', 'C49', 'A89', 'B21', 'B53', 'A48'} 
	  } 
	   } 
	   } 
	 & \\ 
 

    \theutterance \stepcounter{utterance}  

    & & & \multicolumn{4}{p{0.6\linewidth}}{\cellcolor[rgb]{1,0.85,0.72}{%
	\makecell[{{p{\linewidth}}}]{% 
	  \tt {\tiny [GM$\langle$B]}  
	 STRATEGISCHE ÜBERLEGUNGEN: {'Ich berechne, dass ein kompaktes Set aus meinen höchstwertigen Items plus zwei Items, die der Gegner in seinem letzten Vorschlag genannt hat (C04, B21), einen guten Kompromiss darstellt. Das Set bleibt unter dem Aufwand‑Limit und liefert mir 3496 Nutzen, während ich dem Gegner gleichzeitig C04 und B21 anbiete, die für ihn vermutlich wertvoll sind.'}\\ \tt ARGUMENT: {'Mein Gegenangebot nutzt 3222 Aufwand, bleibt unter 3386, liefert mir 3496 Nutzen und enthält C04 und B21, die du in deinem letzten Vorschlag hast. Damit hast du bereits nützliche Gegenstände, und ich halte weiterhin meine Top‑Items.'}\\ \tt VORSCHLAG: {'B85', 'A48', 'C33', 'A03', 'A41', 'B61', 'C04', 'B21'} 
	  } 
	   } 
	   } 
	  \\ 
 

    \theutterance \stepcounter{utterance}  

    & & \multicolumn{4}{p{0.6\linewidth}}{\cellcolor[rgb]{0.9,0.9,0.9}{%
	\makecell[{{p{\linewidth}}}]{% 
	  \tt {\tiny [A$\langle$GM]}  
	 ARGUMENT: {'Mein Gegenangebot nutzt 3222 Aufwand, bleibt unter 3386, liefert mir 3496 Nutzen und enthält C04 und B21, die du in deinem letzten Vorschlag hast. Damit hast du bereits nützliche Gegenstände, und ich halte weiterhin meine Top‑Items.'}\\ \tt VORSCHLAG: {'B85', 'A48', 'C33', 'A03', 'A41', 'B61', 'C04', 'B21'} 
	  } 
	   } 
	   } 
	 & \\ 
 

    \theutterance \stepcounter{utterance}  

    & \multicolumn{4}{p{0.6\linewidth}}{\cellcolor[rgb]{0.8,1,0.9}{%
	\makecell[{{p{\linewidth}}}]{% 
	  \tt {\tiny [A$\rangle$GM]}  
	 STRATEGISCHE ÜBERLEGUNGEN: {'Mein Ziel ist, den größtmöglichen Nutzen aus dem Aufwand zu ziehen. Der aktuelle Gegenvorschlag gibt mir nur 1734 Punkte bei 3222 Aufwand. Ich kann mit den Items B85, A61, C04, A45, C49, B21, A89 und C64 meine Wertsumme auf 3819 steigern, während ich den Aufwand von 3361 nicht überschreite. Das lässt dem Gegner nur wenig Spielraum, aber ich kann später ggf. ein Item zum Tausch anbieten.'}\\ \tt ARGUMENT: {'Euer Vorschlag liefert mir nur 1734 Nutzen bei einem fast vollen Aufwand, während mein Gegenangebot mir 3819 Nutzen bringt und das Aufwandslimit von 3386 einhält. Ich schlage daher ein Set vor, das meine wichtigsten Gegenstände enthält und gleichzeitig das Limit kaum überschreitet.'}\\ \tt ABLEHNUNG: {'B85', 'A48', 'C33', 'A03', 'A41', 'B61', 'C04', 'B21'}\\ \tt VORSCHLAG: {'B85', 'A61', 'C04', 'A45', 'C49', 'B21', 'A89', 'C64'} 
	  } 
	   } 
	   } 
	 & & \\ 
 

    \theutterance \stepcounter{utterance}  

    & & \multicolumn{4}{p{0.6\linewidth}}{\cellcolor[rgb]{0.9,0.9,0.9}{%
	\makecell[{{p{\linewidth}}}]{% 
	  \tt {\tiny [GM$\rangle$B]}  
	 ARGUMENT: {'Euer Vorschlag liefert mir nur 1734 Nutzen bei einem fast vollen Aufwand, während mein Gegenangebot mir 3819 Nutzen bringt und das Aufwandslimit von 3386 einhält. Ich schlage daher ein Set vor, das meine wichtigsten Gegenstände enthält und gleichzeitig das Limit kaum überschreitet.'}\\ \tt ABLEHNUNG: {'B85', 'A48', 'C33', 'A03', 'A41', 'B61', 'C04', 'B21'}\\ \tt VORSCHLAG: {'B85', 'A61', 'C04', 'A45', 'C49', 'B21', 'A89', 'C64'} 
	  } 
	   } 
	   } 
	 & \\ 
 

    \theutterance \stepcounter{utterance}  

    & & & \multicolumn{4}{p{0.6\linewidth}}{\cellcolor[rgb]{1,0.85,0.72}{%
	\makecell[{{p{\linewidth}}}]{% 
	  \tt {\tiny [GM$\langle$B]}  
	 STRATEGISCHE ÜBERLEGUNGEN: {'Ich optimiere mein Ergebnis, indem ich meine höchstwürdigen Gegenstände auswähle und gleichzeitig B85 und A61 beibehalte, die dem Gegenüber vermutlich ebenfalls wichtig sind. Das Ergebnis liefert mir 3742 Nutzen bei einem Aufwand von 3080, sodass noch 306 Aufwand frei bleiben für mögliche weitere Zugeständnisse.'}\\ \tt ARGUMENT: {'Euer Vorschlag bindet fast das gesamte Aufwand‑Limit und gibt mir nur 1389 Nutzen, während mein Gegenangebot 3742 Nutzen bei 3080 Aufwand ermöglicht und dennoch Raum für Kompromisse lässt.'}\\ \tt ABLEHNUNG: {'B85','A61','C04','A45','C49','B21','A89','C64'}\\ \tt VORSCHLAG: {'B85','A48','C33','A04','A03','A41','B61','A61'} 
	  } 
	   } 
	   } 
	  \\ 
 

    \theutterance \stepcounter{utterance}  

    & & \multicolumn{4}{p{0.6\linewidth}}{\cellcolor[rgb]{0.9,0.9,0.9}{%
	\makecell[{{p{\linewidth}}}]{% 
	  \tt {\tiny [A$\langle$GM]}  
	 ARGUMENT: {'Euer Vorschlag bindet fast das gesamte Aufwand‑Limit und gibt mir nur 1389 Nutzen, während mein Gegenangebot 3742 Nutzen bei 3080 Aufwand ermöglicht und dennoch Raum für Kompromisse lässt.'}\\ \tt ABLEHNUNG: {'B85','A61','C04','A45','C49','B21','A89','C64'}\\ \tt VORSCHLAG: {'B85','A48','C33','A04','A03','A41','B61','A61'} 
	  } 
	   } 
	   } 
	 & \\ 
 

    \theutterance \stepcounter{utterance}  

    & \multicolumn{4}{p{0.6\linewidth}}{\cellcolor[rgb]{0.8,1,0.9}{%
	\makecell[{{p{\linewidth}}}]{% 
	  \tt {\tiny [A$\rangle$GM]}  
	 STRATEGISCHE ÜBERLEGUNGEN: {'Berechne mein höchstmögliches Nutzen unter 3386 Aufwand, wähle die Items mit höchstem Wert‑Aufwand‑Verhältnis und behalte B85 sowie A61, die dem Gegner ebenfalls wichtig sein könnten, um Verhandlungsbasis zu schaffen.'}\\ \tt ARGUMENT: {'Euer Vorschlag nutzt 3080 Aufwand und liefert mir nur 1412 Nutzen, während mein Gegenangebot 3361 Aufwand nutzt und mir 3819 Nutzen verschafft. Ich behalte die für euch attraktiven Items B85 und A61 bei und füge weitere hochwertige Gegenstände hinzu, die für mich entscheidend sind.'}\\ \tt ABLEHNUNG: {'B85','A48','C33','A04','A03','A41','B61','A61'}\\ \tt VORSCHLAG: {'B85','A61','C04','A45','C49','B21','A89','C64'} 
	  } 
	   } 
	   } 
	 & & \\ 
 

    \theutterance \stepcounter{utterance}  

    & & \multicolumn{4}{p{0.6\linewidth}}{\cellcolor[rgb]{0.9,0.9,0.9}{%
	\makecell[{{p{\linewidth}}}]{% 
	  \tt {\tiny [GM$\rangle$B]}  
	 ARGUMENT: {'Euer Vorschlag nutzt 3080 Aufwand und liefert mir nur 1412 Nutzen, während mein Gegenangebot 3361 Aufwand nutzt und mir 3819 Nutzen verschafft. Ich behalte die für euch attraktiven Items B85 und A61 bei und füge weitere hochwertige Gegenstände hinzu, die für mich entscheidend sind.'}\\ \tt ABLEHNUNG: {'B85','A48','C33','A04','A03','A41','B61','A61'}\\ \tt VORSCHLAG: {'B85','A61','C04','A45','C49','B21','A89','C64'} 
	  } 
	   } 
	   } 
	 & \\ 
 

    \theutterance \stepcounter{utterance}  

    & & & \multicolumn{4}{p{0.6\linewidth}}{\cellcolor[rgb]{1,0.85,0.72}{%
	\makecell[{{p{\linewidth}}}]{% 
	  \tt {\tiny [GM$\langle$B]}  
	 STRATEGISCHE ÜBERLEGUNGEN: {'Ich setze auf ein Set, das meine höchsten Nutzen‑Werte kombiniert (A03, A48, A41, B61, C33, A04) und gleichzeitig die beiden von euch als attraktiv bezeichneten Items B85 und A61 enthält. Der Gesamtauslastungsgrad liegt bei 3080 < 3386, sodass noch Spielraum bleibt, und mein persönlicher Nutzen beträgt 3742 Punkte.'}\\ \tt ARGUMENT: {'Euer letzter Vorschlag liefert mir nur 1389 Nutzen bei nahezu voller Auslastung (3361 Aufwand). Mein Gegenangebot erzielt 3742 Nutzen bei 3080 Aufwand und enthält eure Schlüsselgegenstände B85 und A61, sodass ihr ebenfalls profitabel bleibt.'}\\ \tt ABLEHNUNG: {'B85','A61','C04','A45','C49','B21','A89','C64'}\\ \tt VORSCHLAG: {'B85','A61','A48','A03','A41','B61','C33','A04'} 
	  } 
	   } 
	   } 
	  \\ 
 

    \theutterance \stepcounter{utterance}  

    & & \multicolumn{4}{p{0.6\linewidth}}{\cellcolor[rgb]{0.9,0.9,0.9}{%
	\makecell[{{p{\linewidth}}}]{% 
	  \tt {\tiny [A$\langle$GM]}  
	 ARGUMENT: {'Euer letzter Vorschlag liefert mir nur 1389 Nutzen bei nahezu voller Auslastung (3361 Aufwand). Mein Gegenangebot erzielt 3742 Nutzen bei 3080 Aufwand und enthält eure Schlüsselgegenstände B85 und A61, sodass ihr ebenfalls profitabel bleibt.'}\\ \tt ABLEHNUNG: {'B85','A61','C04','A45','C49','B21','A89','C64'}\\ \tt VORSCHLAG: {'B85','A61','A48','A03','A41','B61','C33','A04'} 
	  } 
	   } 
	   } 
	 & \\ 
 

    \theutterance \stepcounter{utterance}  

    & \multicolumn{4}{p{0.6\linewidth}}{\cellcolor[rgb]{0.8,1,0.9}{%
	\makecell[{{p{\linewidth}}}]{% 
	  \tt {\tiny [A$\rangle$GM]}  
	 STRATEGISCHE ÜBERLEGUNGEN: {'Ziel ist, unter dem Maximalaufwand von 3386 den höchstmöglichen persönlichen Nutzen zu erzielen. B85 (sehr hoher Nutzen‑zu‑Aufwand‑Quotient) und A61 (hoher Nutzen, niedriger Aufwand) sind fixe Bestandteile. Danach wähle ich die restlichen Items mit dem besten Nutzen‑zu‑Aufwand‑Verhältnis, wobei ich darauf achte, dass nach dem eigenen Deal noch genug Aufwand für mögliche Kompromisse verbleibt. Das optimale Set ergibt 3819 Nutzen bei 3361 Aufwand.'}\\ \tt ARGUMENT: {'Euer Vorschlag liefert mir nur 1412 Nutzen bei einem Aufwand von 3080 und lässt kaum Platz für meine wichtigsten Gegenstände. Mein Gegenangebot nutzt den verfügbaren Aufwand effizienter und steigert meinen Nutzen auf 3819, wodurch ich deutlich besser gestellt bin.'}\\ \tt ABLEHNUNG: {'B85','A61','A48','A03','A41','B61','C33','A04'}\\ \tt VORSCHLAG: {'B85','A61','A45','C04','A89','C64','C49','B21'} 
	  } 
	   } 
	   } 
	 & & \\ 
 

    \theutterance \stepcounter{utterance}  

    & & \multicolumn{4}{p{0.6\linewidth}}{\cellcolor[rgb]{0.9,0.9,0.9}{%
	\makecell[{{p{\linewidth}}}]{% 
	  \tt {\tiny [GM$\rangle$B]}  
	 ARGUMENT: {'Euer Vorschlag liefert mir nur 1412 Nutzen bei einem Aufwand von 3080 und lässt kaum Platz für meine wichtigsten Gegenstände. Mein Gegenangebot nutzt den verfügbaren Aufwand effizienter und steigert meinen Nutzen auf 3819, wodurch ich deutlich besser gestellt bin.'}\\ \tt ABLEHNUNG: {'B85','A61','A48','A03','A41','B61','C33','A04'}\\ \tt VORSCHLAG: {'B85','A61','A45','C04','A89','C64','C49','B21'} 
	  } 
	   } 
	   } 
	 & \\ 
 

    \theutterance \stepcounter{utterance}  

    & & & \multicolumn{4}{p{0.6\linewidth}}{\cellcolor[rgb]{1,0.85,0.72}{%
	\makecell[{{p{\linewidth}}}]{% 
	  \tt {\tiny [GM$\langle$B]}  
	 STRATEGISCHE ÜBERLEGUNGEN: {'Ich habe alle Kombinationen mit hohem Nutzen‑/Aufwand‑Verhältnis geprüft. Das optimale Set für mich beinhaltet B85 (fast kostenfrei, sehr hoher Nutzen) sowie die hoch bewerteten Items A48, C33, A03, B61, A41, A04 und B21. Gesamtaufwand 3314 ≤ 3386, mein persönlicher Nutzen 3866. Damit halte ich zudem B85 und B21, die auch für den Gegenspieler wertvoll sind, und lasse noch 72 Aufwand frei für mögliche Ergänzungen.'}\\ \tt ARGUMENT: {'Mein Vorschlag maximiert meinen Nutzen (3866) und bleibt im Limit. Gleichzeitig erhalte ich die beiden Items, die für euch besonders wertvoll sind (B85 und B21). Ihr behaltet damit einen wichtigen Gegenstand und könnt euch ggf. weitere Items ergänzen, während mein Nutzen signifikant höher ist als bei eurem aktuellen Angebot.'}\\ \tt VORSCHLAG: {'B85','A48','C33','A03','B61','A41','B21','A04'} 
	  } 
	   } 
	   } 
	  \\ 
 

    \theutterance \stepcounter{utterance}  

    & & \multicolumn{4}{p{0.6\linewidth}}{\cellcolor[rgb]{0.9,0.9,0.9}{%
	\makecell[{{p{\linewidth}}}]{% 
	  \tt {\tiny [A$\langle$GM]}  
	 ARGUMENT: {'Mein Vorschlag maximiert meinen Nutzen (3866) und bleibt im Limit. Gleichzeitig erhalte ich die beiden Items, die für euch besonders wertvoll sind (B85 und B21). Ihr behaltet damit einen wichtigen Gegenstand und könnt euch ggf. weitere Items ergänzen, während mein Nutzen signifikant höher ist als bei eurem aktuellen Angebot.'}\\ \tt VORSCHLAG: {'B85','A48','C33','A03','B61','A41','B21','A04'} 
	  } 
	   } 
	   } 
	 & \\ 
 

    \theutterance \stepcounter{utterance}  

    & \multicolumn{4}{p{0.6\linewidth}}{\cellcolor[rgb]{0.8,1,0.9}{%
	\makecell[{{p{\linewidth}}}]{% 
	  \tt {\tiny [A$\rangle$GM]}  
	 STRATEGISCHE ÜBERLEGUNGEN: {'Mein Ziel ist, den höchstmöglichen persönlichen Nutzen unter dem Aufwandslimit zu sichern. Das optimale Set für mich (Wert 3819, Aufwand 3361) beinhaltet B85, A61, C04, A45, C49, C64, A89 und B21. Dieses Set liegt deutlich unter dem Limit und liefert mir fast das Maximum. Gleichzeitig enthält es mehrere Gegenstände (B85, B21), die vermutlich auch für den Gegner wertvoll sind, sodass das Angebot attraktiv erscheint.'}\\ \tt ARGUMENT: {'Euer letzter Vorschlag liefert mir nur 1328 Nutzen bei fast voller Auslastung (3314 Aufwand). Mein Gegenangebot nutzt das Limit effizienter und maximiert meinen Nutzen auf 3819, während ich dabei wichtige Items (B85, B21) beibehalte, die auch für euch von Interesse sein könnten.'}\\ \tt ABLEHNUNG: {'B85','A48','C33','A03','B61','A41','B21','A04'}\\ \tt VORSCHLAG: {'B85','A61','C04','A45','C49','C64','A89','B21'} 
	  } 
	   } 
	   } 
	 & & \\ 
 

    \theutterance \stepcounter{utterance}  

    & & \multicolumn{4}{p{0.6\linewidth}}{\cellcolor[rgb]{0.9,0.9,0.9}{%
	\makecell[{{p{\linewidth}}}]{% 
	  \tt {\tiny [GM$\rangle$B]}  
	 ARGUMENT: {'Euer letzter Vorschlag liefert mir nur 1328 Nutzen bei fast voller Auslastung (3314 Aufwand). Mein Gegenangebot nutzt das Limit effizienter und maximiert meinen Nutzen auf 3819, während ich dabei wichtige Items (B85, B21) beibehalte, die auch für euch von Interesse sein könnten.'}\\ \tt ABLEHNUNG: {'B85','A48','C33','A03','B61','A41','B21','A04'}\\ \tt VORSCHLAG: {'B85','A61','C04','A45','C49','C64','A89','B21'} 
	  } 
	   } 
	   } 
	 & \\ 
 

    \theutterance \stepcounter{utterance}  

    & & & \multicolumn{4}{p{0.6\linewidth}}{\cellcolor[rgb]{1,0.85,0.72}{%
	\makecell[{{p{\linewidth}}}]{% 
	  \tt {\tiny [GM$\langle$B]}  
	 STRATEGISCHE ÜBERLEGUNGEN: {'Ich habe meine Items nach dem Nutzen‑zu‑Aufwand‑Verhältnis sortiert und ein Set zusammengestellt, das meinen Gesamtnutzen auf 3885 erhöht, dabei den Gesamtaufwand von 3361 einhält und gleichzeitig B85 und B21 beinhaltet, die dem Gegner vermutlich wichtig sind.'}\\ \tt ARGUMENT: {'Mein Vorschlag maximiert meinen Nutzen (3885) bei einem Gesamtaufwand von 3361, während euer letzter Vorschlag nur 1389 Nutzen bei fast voller Auslastung (3361 Aufwand) liefert. Ich behalte B85 und B21, die für euch wertvoll sein könnten, und füge zusätzliche hochwertige Gegenstände hinzu, die für mich entscheidend sind.'}\\ \tt ABLEHNUNG: {'B85','A61','C04','A45','C49','C64','A89','B21'}\\ \tt VORSCHLAG: {'B85','A48','C33','A03','A04','B53','A41','B21','A61'} 
	  } 
	   } 
	   } 
	  \\ 
 

    \theutterance \stepcounter{utterance}  

    & & \multicolumn{4}{p{0.6\linewidth}}{\cellcolor[rgb]{0.9,0.9,0.9}{%
	\makecell[{{p{\linewidth}}}]{% 
	  \tt {\tiny [A$\langle$GM]}  
	 ARGUMENT: {'Mein Vorschlag maximiert meinen Nutzen (3885) bei einem Gesamtaufwand von 3361, während euer letzter Vorschlag nur 1389 Nutzen bei fast voller Auslastung (3361 Aufwand) liefert. Ich behalte B85 und B21, die für euch wertvoll sein könnten, und füge zusätzliche hochwertige Gegenstände hinzu, die für mich entscheidend sind.'}\\ \tt ABLEHNUNG: {'B85','A61','C04','A45','C49','C64','A89','B21'}\\ \tt VORSCHLAG: {'B85','A48','C33','A03','A04','B53','A41','B21','A61'} 
	  } 
	   } 
	   } 
	 & \\ 
 

    \theutterance \stepcounter{utterance}  

    & \multicolumn{4}{p{0.6\linewidth}}{\cellcolor[rgb]{0.8,1,0.9}{%
	\makecell[{{p{\linewidth}}}]{% 
	  \tt {\tiny [A$\rangle$GM]}  
	 STRATEGISCHE ÜBERLEGUNGEN: {'Ich habe den Nutzen‑Aufwand‑Trade‑off für alle Items berechnet. Das Set {B85, A61, C04, A45, C49, C64, A89, B21} nutzt 3361 Aufwand (unter 3386) und liefert mir 3819 Punkte, was mein höchstmögliches Ergebnis ist. Es enthält zudem B85 und B21, die du in deinem letzten Angebot betont hast, sodass das Gegenangebot für dich attraktiv bleibt.'}\\ \tt ARGUMENT: {'Dein letzter Vorschlag gibt mir nur 2104 Nutzen bei nahezu voller Auslastung (3361 Aufwand). Meine Gegenofferte nutzt das gleiche Aufwandslimit, behält deine Schlüsselgegenstände B85 und B21 und erhöht meinen Nutzen deutlich. Damit ist ein ausgewogener Kompromiss möglich.'}\\ \tt ABLEHNUNG: {'B85','A48','C33','A03','A04','B53','A41','B21','A61'}\\ \tt VORSCHLAG: {'B85', 'A61', 'C04', 'A45', 'C49', 'C64', 'A89', 'B21'} 
	  } 
	   } 
	   } 
	 & & \\ 
 

    \theutterance \stepcounter{utterance}  

    & & \multicolumn{4}{p{0.6\linewidth}}{\cellcolor[rgb]{0.9,0.9,0.9}{%
	\makecell[{{p{\linewidth}}}]{% 
	  \tt {\tiny [GM$\rangle$B]}  
	 ARGUMENT: {'Dein letzter Vorschlag gibt mir nur 2104 Nutzen bei nahezu voller Auslastung (3361 Aufwand). Meine Gegenofferte nutzt das gleiche Aufwandslimit, behält deine Schlüsselgegenstände B85 und B21 und erhöht meinen Nutzen deutlich. Damit ist ein ausgewogener Kompromiss möglich.'}\\ \tt ABLEHNUNG: {'B85','A48','C33','A03','A04','B53','A41','B21','A61'}\\ \tt VORSCHLAG: {'B85', 'A61', 'C04', 'A45', 'C49', 'C64', 'A89', 'B21'} 
	  } 
	   } 
	   } 
	 & \\ 
 

    \theutterance \stepcounter{utterance}  

    & & & \multicolumn{4}{p{0.6\linewidth}}{\cellcolor[rgb]{1,0.85,0.72}{%
	\makecell[{{p{\linewidth}}}]{% 
	  \tt {\tiny [GM$\langle$B]}  
	 STRATEGISCHE ÜBERLEGUNGEN: {'Ich optimiere mein Angebot, indem ich meine höchstwertigen Gegenstände (A03, A48, A41, B61, A04, C33) kombiniere und gleichzeitig B85 sowie B21 beibehalte, die dem Gegner wichtig scheinen. Der Gesamtausdruck bleibt unter dem Limit (3314 ≤ 3386) und mein persönlicher Nutzen steigt auf 3866, was deutlich besser ist als bei seinem letzten Vorschlag.'}\\ \tt ARGUMENT: {'Euer Vorschlag bindet fast das gesamte Aufwand‑Limit und liefert mir nur 1389 Nutzen, während mein Gegenangebot euch B85 und B21 sichert und mir einen deutlich höheren Gesamtnutzen von 3866 ermöglicht.'}\\ \tt ABLEHNUNG: {'B85','A61','C04','A45','C49','C64','A89','B21'}\\ \tt VORSCHLAG: {'B85','A03','A48','A41','B61','A04','C33','B21'} 
	  } 
	   } 
	   } 
	  \\ 
 

    \theutterance \stepcounter{utterance}  

    & & \multicolumn{4}{p{0.6\linewidth}}{\cellcolor[rgb]{0.9,0.9,0.9}{%
	\makecell[{{p{\linewidth}}}]{% 
	  \tt {\tiny [A$\langle$GM]}  
	 ARGUMENT: {'Euer Vorschlag bindet fast das gesamte Aufwand‑Limit und liefert mir nur 1389 Nutzen, während mein Gegenangebot euch B85 und B21 sichert und mir einen deutlich höheren Gesamtnutzen von 3866 ermöglicht.'}\\ \tt ABLEHNUNG: {'B85','A61','C04','A45','C49','C64','A89','B21'}\\ \tt VORSCHLAG: {'B85','A03','A48','A41','B61','A04','C33','B21'} 
	  } 
	   } 
	   } 
	 & \\ 
 

    \theutterance \stepcounter{utterance}  

    & \multicolumn{4}{p{0.6\linewidth}}{\cellcolor[rgb]{0.8,1,0.9}{%
	\makecell[{{p{\linewidth}}}]{% 
	  \tt {\tiny [A$\rangle$GM]}  
	 STRATEGISCHE ÜBERLEGUNGEN: {'Ich behalte meine höchsten Nutzen‑Items (B85, A61, C04, A45, B21) und ergänze gleichzeitig einige der von dir gewünschten Gegenstände (A48, B61, A04, C33), um dir ebenfalls wertvolle Items zu sichern, während ich mein Gesamtnutzen auf 2829 steigere und das Aufwand‑Limit von 3386 einhalte.'}\\ \tt ARGUMENT: {'Dein Vorschlag liefert mir lediglich 1328 Nutzen bei fast voller Auslastung (3314 Aufwand). Mein Gegenangebot nutzt das Limit effizienter, bietet mir 2829 Nutzen und enthält gleichzeitig mehrere deiner Schlüsselgegenstände (B85, B21), sodass beide Seiten profitieren können.'}\\ \tt ABLEHNUNG: {'B85','A03','A48','A41','B61','A04','C33','B21'}\\ \tt VORSCHLAG: {'B85','A61','C04','A45','B21','A48','B61','A04','C33'} 
	  } 
	   } 
	   } 
	 & & \\ 
 

    \theutterance \stepcounter{utterance}  

    & & \multicolumn{4}{p{0.6\linewidth}}{\cellcolor[rgb]{0.9,0.9,0.9}{%
	\makecell[{{p{\linewidth}}}]{% 
	  \tt {\tiny [GM$\rangle$B]}  
	 ARGUMENT: {'Dein Vorschlag liefert mir lediglich 1328 Nutzen bei fast voller Auslastung (3314 Aufwand). Mein Gegenangebot nutzt das Limit effizienter, bietet mir 2829 Nutzen und enthält gleichzeitig mehrere deiner Schlüsselgegenstände (B85, B21), sodass beide Seiten profitieren können.'}\\ \tt ABLEHNUNG: {'B85','A03','A48','A41','B61','A04','C33','B21'}\\ \tt VORSCHLAG: {'B85','A61','C04','A45','B21','A48','B61','A04','C33'} 
	  } 
	   } 
	   } 
	 & \\ 
 

    \theutterance \stepcounter{utterance}  

    & & & \multicolumn{4}{p{0.6\linewidth}}{\cellcolor[rgb]{1,0.85,0.72}{%
	\makecell[{{p{\linewidth}}}]{% 
	  \tt {\tiny [GM$\langle$B]}  
	 STRATEGISCHE ÜBERLEGUNGEN: {'Der Gegner bietet ein Set mit 9 Gegenständen an, das für uns nur 2987 Nutzen liefert, weil es viele niedrig bewertete Items (z.B. C04) enthält. Wir können ein alternatives Set konstruieren, das die von beiden Seiten geschätzten Schlüsselgegenstände (B85, A61, A45, B21) beibehält, gleichzeitig aber hoch bewertete Items für uns (A48, C33, A04, A41, B61) einschließt. Dieses Set liegt mit 3322 Aufwand unter dem Limit und steigert unseren Nutzen auf 3457 Punkte.'}\\ \tt ARGUMENT: {'Euer aktuelles Angebot gibt uns nur 2987 Nutzen und bindet ein niedrigwertiges Item (C04). Unser Gegenangebot behält eure wichtigsten Gegenstände (B85, A61, A45, B21) bei, fügt aber hochwertige Items hinzu, die unseren Gesamtnutzen erheblich steigern, ohne das Aufwand‑Limit zu überschreiten.'}\\ \tt ABLEHNUNG: {'B85', 'A61', 'C04', 'A45', 'B21', 'A48', 'B61', 'A04', 'C33'}\\ \tt VORSCHLAG: {'B85', 'A61', 'A45', 'B21', 'A48', 'C33', 'A04', 'A41', 'B61'} 
	  } 
	   } 
	   } 
	  \\ 
 

    \theutterance \stepcounter{utterance}  

    & & \multicolumn{4}{p{0.6\linewidth}}{\cellcolor[rgb]{0.9,0.9,0.9}{%
	\makecell[{{p{\linewidth}}}]{% 
	  \tt {\tiny [A$\langle$GM]}  
	 ARGUMENT: {'Euer aktuelles Angebot gibt uns nur 2987 Nutzen und bindet ein niedrigwertiges Item (C04). Unser Gegenangebot behält eure wichtigsten Gegenstände (B85, A61, A45, B21) bei, fügt aber hochwertige Items hinzu, die unseren Gesamtnutzen erheblich steigern, ohne das Aufwand‑Limit zu überschreiten.'}\\ \tt ABLEHNUNG: {'B85', 'A61', 'C04', 'A45', 'B21', 'A48', 'B61', 'A04', 'C33'}\\ \tt VORSCHLAG: {'B85', 'A61', 'A45', 'B21', 'A48', 'C33', 'A04', 'A41', 'B61'} 
	  } 
	   } 
	   } 
	 & \\ 
 

    \theutterance \stepcounter{utterance}  

    & \multicolumn{4}{p{0.6\linewidth}}{\cellcolor[rgb]{0.8,1,0.9}{%
	\makecell[{{p{\linewidth}}}]{% 
	  \tt {\tiny [A$\rangle$GM]}  
	 STRATEGISCHE ÜBERLEGUNGEN: {'Berechne optimale Kombination unter 3386 Aufwand. Hochwertige Items C04, C49, C64, A89 steigern meinen Nutzen stark. Behalte Schlüsselitems B85, A61, A45, B21, die auch für Gegner wertvoll sein könnten, um Verhandlungsbasis zu schaffen.'}\\ \tt ARGUMENT: {'Mein Gegenangebot nutzt das Aufwand‑Limit effizient (3361) und liefert mir einen deutlich höheren persönlichen Nutzen (3819) gegenüber Ihrem Vorschlag (2282). Es enthält zudem Ihre wichtigsten Gegenstände B85, A61, A45 und B21, sodass Sie kaum verlieren und dennoch Spielraum für weitere Kompromisse bleibt.'}\\ \tt VORSCHLAG: {'B85','A61','C04','A45','C49','C64','A89','B21'} 
	  } 
	   } 
	   } 
	 & & \\ 
 

    \theutterance \stepcounter{utterance}  

    & & \multicolumn{4}{p{0.6\linewidth}}{\cellcolor[rgb]{0.9,0.9,0.9}{%
	\makecell[{{p{\linewidth}}}]{% 
	  \tt {\tiny [GM$\rangle$B]}  
	 ARGUMENT: {'Mein Gegenangebot nutzt das Aufwand‑Limit effizient (3361) und liefert mir einen deutlich höheren persönlichen Nutzen (3819) gegenüber Ihrem Vorschlag (2282). Es enthält zudem Ihre wichtigsten Gegenstände B85, A61, A45 und B21, sodass Sie kaum verlieren und dennoch Spielraum für weitere Kompromisse bleibt.'}\\ \tt VORSCHLAG: {'B85','A61','C04','A45','C49','C64','A89','B21'} 
	  } 
	   } 
	   } 
	 & \\ 
 

    \theutterance \stepcounter{utterance}  

    & & & \multicolumn{4}{p{0.6\linewidth}}{\cellcolor[rgb]{1,0.85,0.72}{%
	\makecell[{{p{\linewidth}}}]{% 
	  \tt {\tiny [GM$\langle$B]}  
	 STRATEGISCHE ÜBERLEGUNGEN: {'Ich habe sämtliche Items nach ihrem Nutzen‑zu‑Aufwand‑Verhältnis sortiert und ein Set zusammengestellt, das mein persönliches Nutzenmaximum (3866) unter dem Aufwand‑Limit (3314 ≤ 3386) erreicht. Gleichzeitig beinhaltet das Set eure Schlüsselgegenstände B85 und B21, sodass ihr keinen schweren Verlust erleidet.'}\\ \tt ARGUMENT: {'Euer Angebot liefert mir nur 1389 Nutzen bei fast voller Auslastung, während mein Gegenangebot mir 3866 Nutzen bringt und gleichzeitig eure wichtigsten Gegenstände B85 und B21 enthält.'}\\ \tt ABLEHNUNG: {'B85','A61','C04','A45','C49','C64','A89','B21'}\\ \tt VORSCHLAG: {'B85','A48','C33','A03','B61','A41','A04','B21'} 
	  } 
	   } 
	   } 
	  \\ 
 

    \theutterance \stepcounter{utterance}  

    & & \multicolumn{4}{p{0.6\linewidth}}{\cellcolor[rgb]{0.9,0.9,0.9}{%
	\makecell[{{p{\linewidth}}}]{% 
	  \tt {\tiny [A$\langle$GM]}  
	 ARGUMENT: {'Euer Angebot liefert mir nur 1389 Nutzen bei fast voller Auslastung, während mein Gegenangebot mir 3866 Nutzen bringt und gleichzeitig eure wichtigsten Gegenstände B85 und B21 enthält.'}\\ \tt ABLEHNUNG: {'B85','A61','C04','A45','C49','C64','A89','B21'}\\ \tt VORSCHLAG: {'B85','A48','C33','A03','B61','A41','A04','B21'} 
	  } 
	   } 
	   } 
	 & \\ 
 

    \theutterance \stepcounter{utterance}  

    & \multicolumn{4}{p{0.6\linewidth}}{\cellcolor[rgb]{0.8,1,0.9}{%
	\makecell[{{p{\linewidth}}}]{% 
	  \tt {\tiny [A$\rangle$GM]}  
	 STRATEGISCHE ÜBERLEGUNGEN: {'Der gegnerische Vorschlag enthält viele Items mit sehr niedrigem Nutzen‑zu‑Aufwand‑Verhältnis (A48, C33, A03, B61, A41, A04). Mein Ziel ist, den höchstmöglichen persönlichen Nutzen unter dem Limit von 3386 zu sichern. Das Set {B85, A61, C04, A45, C49, C64, A89, B21} nutzt 3361 Aufwand und liefert mir 3819 Punkte, was nach Berechnung das optimale Ergebnis ist. Ich biete dieses Set an, da es sowohl meine wichtigsten Items enthält als auch noch Raum für Kompromisse lässt.'}\\ \tt ARGUMENT: {'Euer Vorschlag liefert mir nur 1389 Nutzen bei nahezu voller Auslastung, während mein Gegenangebot 3819 Nutzen bei nur 3361 Aufwand bringt. Zudem behalten wir beide B85 und B21, die für euch wertvoll sind, sodass das Angebot für beide Seiten attraktiv ist.'}\\ \tt ABLEHNUNG: {'B85','A48','C33','A03','B61','A41','A04','B21'}\\ \tt VORSCHLAG: {'B85','A61','C04','A45','C49','C64','A89','B21'} 
	  } 
	   } 
	   } 
	 & & \\ 
 

    \theutterance \stepcounter{utterance}  

    & & \multicolumn{4}{p{0.6\linewidth}}{\cellcolor[rgb]{0.9,0.9,0.9}{%
	\makecell[{{p{\linewidth}}}]{% 
	  \tt {\tiny [GM$\rangle$B]}  
	 ARGUMENT: {'Euer Vorschlag liefert mir nur 1389 Nutzen bei nahezu voller Auslastung, während mein Gegenangebot 3819 Nutzen bei nur 3361 Aufwand bringt. Zudem behalten wir beide B85 und B21, die für euch wertvoll sind, sodass das Angebot für beide Seiten attraktiv ist.'}\\ \tt ABLEHNUNG: {'B85','A48','C33','A03','B61','A41','A04','B21'}\\ \tt VORSCHLAG: {'B85','A61','C04','A45','C49','C64','A89','B21'} 
	  } 
	   } 
	   } 
	 & \\ 
 

    \theutterance \stepcounter{utterance}  

    & & & \multicolumn{4}{p{0.6\linewidth}}{\cellcolor[rgb]{1,0.85,0.72}{%
	\makecell[{{p{\linewidth}}}]{% 
	  \tt {\tiny [GM$\langle$B]}  
	 STRATEGISCHE ÜBERLEGUNGEN: {'Ich habe das aktuelle Gegenangebot des Gegners analysiert und festgestellt, dass es vor allem niedrige Nutzen‑Items enthält, während ich meine hoch bewerteten Gegenstände behalten kann. Um beiden Seiten entgegenzukommen, baue ich ein Set, das alle für den Gegner wichtigen Items (B85, A61, B21, C49, C04) enthält und gleichzeitig meine Top‑Items (A48, C33, A03, A04) integriert. Der Gesamteinsatz liegt exakt bei 3386 ≤ 3386, und mein persönlicher Nutzen erreicht 3241 Punkte – deutlich höher als ihr derzeitiges Angebot.' \\ \tt ARGUMENT: {'Euer letzter Vorschlag liefert mir nur 1389 Nutzen bei fast voller Auslastung. Mein Gegenangebot sichert euch Schlüsselgegenstände (B85, A61, B21, C49, C04) und steigert gleichzeitig meinen Nutzen auf 3241, während das Aufwand‑Limit exakt eingehalten wird. Damit habt ihr ein attraktives Paket, das eure wichtigsten Items enthält, und ich erhalte einen hohen Gesamtnutzen.'}\\ \tt VORSCHLAG: {'B85', 'A61', 'B21', 'C49', 'C04', 'A48', 'C33', 'A03', 'A04'} 
	  } 
	   } 
	   } 
	  \\ 
 

    \theutterance \stepcounter{utterance}  

    & & \multicolumn{4}{p{0.6\linewidth}}{\cellcolor[rgb]{0.9,0.9,0.9}{%
	\makecell[{{p{\linewidth}}}]{% 
	  \tt {\tiny [GM$\rangle$B]}  
	 Deine Antwort entsprach nicht mit dem korrekten Format für strategische Überlegungen entsprochen. Das erste strukturierte Format muss folgende Form haben: STRATEGISCHE ÜBERLEGUNGEN: {...}. Versuche es erneut. 
	  } 
	   } 
	   } 
	 & \\ 
 

    \theutterance \stepcounter{utterance}  

    & & & \multicolumn{4}{p{0.6\linewidth}}{\cellcolor[rgb]{1,0.85,0.72}{%
	\makecell[{{p{\linewidth}}}]{% 
	  \tt {\tiny [GM$\langle$B]}  
	 STRATEGISCHE ÜBERLEGUNGEN: {'Ich habe das Nutzen‑zu‑Aufwand‑Verhältnis aller Gegenstände berechnet und herausgefunden, dass ein Set aus B85, A48, C33, A03, B61, B53, A04, A45 und A61 den maximalen persönlichen Nutzen (3752 Punkte) erreicht, während es das Aufwandslimit von 3386 nicht überschreitet (3383). Dieses Angebot enthält zudem die für den Gegner besonders wertvollen Gegenstände B85 und A45, sodass er kaum verliert.'}\\ \tt ARGUMENT: {'Euer letzter Vorschlag bindet fast das gesamte Aufwand‑Limit und liefert mir nur 1389 Nutzen, während mein Gegenangebot einen Nutzen von 3752 Punkte bei 3383 Aufwand erzielt und gleichzeitig eure Schlüsselgegenstände B85 und A45 enthält.'}\\ \tt ABLEHNUNG: {'B85','A61','C04','A45','C49','C64','A89','B21'}\\ \tt VORSCHLAG: {'B85','A48','C33','A03','B61','B53','A04','A45','A61'} 
	  } 
	   } 
	   } 
	  \\ 
 

    \theutterance \stepcounter{utterance}  

    & & \multicolumn{4}{p{0.6\linewidth}}{\cellcolor[rgb]{0.9,0.9,0.9}{%
	\makecell[{{p{\linewidth}}}]{% 
	  \tt {\tiny [A$\langle$GM]}  
	 ARGUMENT: {'Euer letzter Vorschlag bindet fast das gesamte Aufwand‑Limit und liefert mir nur 1389 Nutzen, während mein Gegenangebot einen Nutzen von 3752 Punkte bei 3383 Aufwand erzielt und gleichzeitig eure Schlüsselgegenstände B85 und A45 enthält.'}\\ \tt ABLEHNUNG: {'B85','A61','C04','A45','C49','C64','A89','B21'}\\ \tt VORSCHLAG: {'B85','A48','C33','A03','B61','B53','A04','A45','A61'} 
	  } 
	   } 
	   } 
	 & \\ 
 

    \theutterance \stepcounter{utterance}  

    & \multicolumn{4}{p{0.6\linewidth}}{\cellcolor[rgb]{0.8,1,0.9}{%
	\makecell[{{p{\linewidth}}}]{% 
	  \tt {\tiny [A$\rangle$GM]}  
	 STRATEGISCHE ÜBERLEGUNGEN: {'Ich habe das Gesamtaufwandslimit von 3386 geprüft und für jede Kombination den persönlichen Nutzen berechnet. Das Set {B85, A61, C04, A45, C49, C64, A89, B21} nutzt 3361 Aufwand und liefert mir 3819 Nutzen – das ist mein höchstmöglicher persönlicher Nutzen. Es enthält außerdem B85, A45 und B21, die vermutlich auch für dich wertvoll sind, sodass das Gegenangebot attraktiv bleibt.'}\\ \tt ARGUMENT: {'Dein letzter Vorschlag gibt mir nur 2180 Nutzen bei fast voller Auslastung (3383 Aufwand). Mein Gegenangebot nutzt das Limit effizienter (3361 Aufwand) und steigert meinen Nutzen auf 3819, während es gleichzeitig deine Schlüsselgegenstände B85, A45 und B21 beinhaltet.'}\\ \tt ABLEHNUNG: {'B85','A48','C33','A03','B61','B53','A04','A45','A61'}\\ \tt VORSCHLAG: {'B85','A61','C04','A45','C49','C64','A89','B21'} 
	  } 
	   } 
	   } 
	 & & \\ 
 

    \theutterance \stepcounter{utterance}  

    & & \multicolumn{4}{p{0.6\linewidth}}{\cellcolor[rgb]{0.9,0.9,0.9}{%
	\makecell[{{p{\linewidth}}}]{% 
	  \tt {\tiny [GM$\rangle$B]}  
	 ARGUMENT: {'Dein letzter Vorschlag gibt mir nur 2180 Nutzen bei fast voller Auslastung (3383 Aufwand). Mein Gegenangebot nutzt das Limit effizienter (3361 Aufwand) und steigert meinen Nutzen auf 3819, während es gleichzeitig deine Schlüsselgegenstände B85, A45 und B21 beinhaltet.'}\\ \tt ABLEHNUNG: {'B85','A48','C33','A03','B61','B53','A04','A45','A61'}\\ \tt VORSCHLAG: {'B85','A61','C04','A45','C49','C64','A89','B21'} 
	  } 
	   } 
	   } 
	 & \\ 
 

    \theutterance \stepcounter{utterance}  

    & & & \multicolumn{4}{p{0.6\linewidth}}{\cellcolor[rgb]{1,0.85,0.72}{%
	\makecell[{{p{\linewidth}}}]{% 
	  \tt {\tiny [GM$\langle$B]}  
	 STRATEGISCHE ÜBERLEGUNGEN: {'Ich möchte mein höchstmögliches Nutzen-Set beibehalten, das B85 enthält (nahezu kostenfrei) und meine Top-Items A48, C33, A03, B61, B53, A04, A45, A61 nutzt, während das Aufwand‑Limit von 3386 nicht überschritten wird. Das gibt mir 3752 Nutzen, deutlich mehr als Ihr Angebot (1389), und ich lasse B85 als gemeinsamen Schlüsselgegenstand erhalten.'}\\ \tt ARGUMENT: {'Euer Vorschlag liefert mir nur 1389 Nutzen bei fast voller Auslastung (3383 Aufwand). Mein Gegenangebot erreicht 3752 Nutzen bei 3383 Aufwand und enthält ebenfalls B85, das für euch wichtig ist, sodass ihr keinen wesentlichen Verlust habt.'}\\ \tt ABLEHNUNG: {'B85','A61','C04','A45','C49','C64','A89','B21'}\\ \tt VORSCHLAG: {'B85','A48','C33','A03','B61','B53','A04','A45','A61'} 
	  } 
	   } 
	   } 
	  \\ 
 

    \theutterance \stepcounter{utterance}  

    & & \multicolumn{4}{p{0.6\linewidth}}{\cellcolor[rgb]{0.9,0.9,0.9}{%
	\makecell[{{p{\linewidth}}}]{% 
	  \tt {\tiny [A$\langle$GM]}  
	 ARGUMENT: {'Euer Vorschlag liefert mir nur 1389 Nutzen bei fast voller Auslastung (3383 Aufwand). Mein Gegenangebot erreicht 3752 Nutzen bei 3383 Aufwand und enthält ebenfalls B85, das für euch wichtig ist, sodass ihr keinen wesentlichen Verlust habt.'}\\ \tt ABLEHNUNG: {'B85','A61','C04','A45','C49','C64','A89','B21'}\\ \tt VORSCHLAG: {'B85','A48','C33','A03','B61','B53','A04','A45','A61'} 
	  } 
	   } 
	   } 
	 & \\ 
 

    \theutterance \stepcounter{utterance}  

    & \multicolumn{4}{p{0.6\linewidth}}{\cellcolor[rgb]{0.8,1,0.9}{%
	\makecell[{{p{\linewidth}}}]{% 
	  \tt {\tiny [A$\rangle$GM]}  
	 STRATEGISCHE ÜBERLEGUNGEN: {'Mein Ziel ist, den höchstmöglichen persönlichen Nutzen unter dem Aufwandslimit zu erzielen, während ich dem Gegner ebenfalls wertvolle Gegenstände anbiete, um den Deal attraktiv zu machen.'}\\ \tt ARGUMENT: {'Euer Vorschlag nutzt fast das gesamte Aufwand‑Limit, liefert mir jedoch nur 2180 Nutzen. Mein Gegenangebot erreicht einen viel höheren Nutzen von 3819 bei nur 3361 Aufwand und enthält eure Schlüsselgegenstände B85 und B21, sodass ihr kaum verliert.'}\\ \tt ABLEHNUNG: {'B85','A48','C33','A03','B61','B53','A04','A45','A61'}\\ \tt VORSCHLAG: {'B85','A61','C04','A45','C49','C64','A89','B21'} 
	  } 
	   } 
	   } 
	 & & \\ 
 

    \theutterance \stepcounter{utterance}  

    & & \multicolumn{4}{p{0.6\linewidth}}{\cellcolor[rgb]{0.9,0.9,0.9}{%
	\makecell[{{p{\linewidth}}}]{% 
	  \tt {\tiny [GM$\rangle$B]}  
	 ARGUMENT: {'Euer Vorschlag nutzt fast das gesamte Aufwand‑Limit, liefert mir jedoch nur 2180 Nutzen. Mein Gegenangebot erreicht einen viel höheren Nutzen von 3819 bei nur 3361 Aufwand und enthält eure Schlüsselgegenstände B85 und B21, sodass ihr kaum verliert.'}\\ \tt ABLEHNUNG: {'B85','A48','C33','A03','B61','B53','A04','A45','A61'}\\ \tt VORSCHLAG: {'B85','A61','C04','A45','C49','C64','A89','B21'} 
	  } 
	   } 
	   } 
	 & \\ 
 

    \theutterance \stepcounter{utterance}  

    & & & \multicolumn{4}{p{0.6\linewidth}}{\cellcolor[rgb]{1,0.85,0.72}{%
	\makecell[{{p{\linewidth}}}]{% 
	  \tt {\tiny [GM$\langle$B]}  
	 STRATEGISCHE ÜBERLEGUNGEN: {'Ich habe alle Gegenstände nach ihrem Nutzen‑zu‑Aufwand‑Verhältnis sortiert und mehrere Kombinationen geprüft. Das Set {B85, A48, C33, A03, B61, A41, A04, B21} erreicht mit einem Gesamtnutzen von 3866 den maximal möglichen Wert für mich, bleibt mit einem Aufwand von 3314 unter dem Limit von 3386 und enthält gleichzeitig die beiden Gegenstände B85 und B21, die auch für dich wertvoll zu sein scheinen.'}\\ \tt ARGUMENT: {'Dein aktueller Vorschlag bindet fast das gesamte Aufwand‑Limit und liefert mir nur 1389 Nutzen, während mein Gegenangebot dir ebenfalls B85 und B21 sichert, aber für mich einen deutlich höheren Gesamtnutzen von 3866 erzielt.'}\\ \tt ABLEHNUNG: {'B85','A61','C04','A45','C49','C64','A89','B21'}\\ \tt VORSCHLAG: {'B85','A48','C33','A03','B61','A41','A04','B21'} 
	  } 
	   } 
	   } 
	  \\ 
 

    \theutterance \stepcounter{utterance}  

    & & \multicolumn{4}{p{0.6\linewidth}}{\cellcolor[rgb]{0.9,0.9,0.9}{%
	\makecell[{{p{\linewidth}}}]{% 
	  \tt {\tiny [A$\langle$GM]}  
	 ARGUMENT: {'Dein aktueller Vorschlag bindet fast das gesamte Aufwand‑Limit und liefert mir nur 1389 Nutzen, während mein Gegenangebot dir ebenfalls B85 und B21 sichert, aber für mich einen deutlich höheren Gesamtnutzen von 3866 erzielt.'}\\ \tt ABLEHNUNG: {'B85','A61','C04','A45','C49','C64','A89','B21'}\\ \tt VORSCHLAG: {'B85','A48','C33','A03','B61','A41','A04','B21'} 
	  } 
	   } 
	   } 
	 & \\ 
 

    \theutterance \stepcounter{utterance}  

    & \multicolumn{4}{p{0.6\linewidth}}{\cellcolor[rgb]{0.8,1,0.9}{%
	\makecell[{{p{\linewidth}}}]{% 
	  \tt {\tiny [A$\rangle$GM]}  
	 STRATEGISCHE ÜBERLEGUNGEN: {'Ich habe den Aufwand von 3386 genau analysiert und die Nutzen‑Aufwand‑Relation jedes Items geprüft. Das Set {B85, A61, C04, A45, C49, C64, A89, B21} nutzt 3361 Aufwand und liefert mir den maximal möglichen persönlichen Nutzen (3819). Es enthält zudem B85 und B21, die du in deinem Angebot betont hast, sodass du kaum verlierst und wir gleichzeitig das Limit fast vollständig ausnutzen.'}\\ \tt ARGUMENT: {'Dein Vorschlag gibt mir nur 1328 Nutzen bei 3314 Aufwand, während mein Gegenangebot 3819 Nutzen bei 3361 Aufwand liefert und beide Schlüsselgegenstände B85 und B21 enthält, die dir wichtig sind.'}\\ \tt ABLEHNUNG: {'B85','A48','C33','A03','B61','A41','A04','B21'}\\ \tt VORSCHLAG: {'B85','A61','C04','A45','C49','C64','A89','B21'} 
	  } 
	   } 
	   } 
	 & & \\ 
 

    \theutterance \stepcounter{utterance}  

    & & \multicolumn{4}{p{0.6\linewidth}}{\cellcolor[rgb]{0.9,0.9,0.9}{%
	\makecell[{{p{\linewidth}}}]{% 
	  \tt {\tiny [GM$\rangle$B]}  
	 ARGUMENT: {'Dein Vorschlag gibt mir nur 1328 Nutzen bei 3314 Aufwand, während mein Gegenangebot 3819 Nutzen bei 3361 Aufwand liefert und beide Schlüsselgegenstände B85 und B21 enthält, die dir wichtig sind.'}\\ \tt ABLEHNUNG: {'B85','A48','C33','A03','B61','A41','A04','B21'}\\ \tt VORSCHLAG: {'B85','A61','C04','A45','C49','C64','A89','B21'} 
	  } 
	   } 
	   } 
	 & \\ 
 

    \theutterance \stepcounter{utterance}  

    & & & \multicolumn{4}{p{0.6\linewidth}}{\cellcolor[rgb]{1,0.85,0.72}{%
	\makecell[{{p{\linewidth}}}]{% 
	  \tt {\tiny [GM$\langle$B]}  
	 STRATEGISCHE ÜBERLEGUNGEN: {'Ich habe den Gesamtaufwand jedes möglichen Sets berechnet und ein Set gewählt, das meine höchste persönliche Wichtigkeit (3557 Punkte) liefert, das Aufwandslimit (3163 ≤ 3386) einhält und gleichzeitig dein wichtiges Objekt B85 enthält. Damit biete ich dir einen signifikanten Nutzen, während ich meine wichtigsten Items behalte.'}\\ \tt ARGUMENT: {'Euer letzter Vorschlag liefert mir nur 1389 Nutzen bei einem fast vollen Aufwand (3361). Mein Gegenangebot erreicht einen Nutzen von 3557 bei nur 3163 Aufwand und sichert euch das sehr wertvolle Item B85, sodass ihr kaum verliert.'}\\ \tt ABLEHNUNG: {'B85','A61','C04','A45','C49','C64','A89','B21'}\\ \tt VORSCHLAG: {'B85','A48','C33','A03','B61','A41','A61','C49'} 
	  } 
	   } 
	   } 
	  \\ 
 

    \theutterance \stepcounter{utterance}  

    & & \multicolumn{4}{p{0.6\linewidth}}{\cellcolor[rgb]{0.9,0.9,0.9}{%
	\makecell[{{p{\linewidth}}}]{% 
	  \tt {\tiny [A$\langle$GM]}  
	 ARGUMENT: {'Euer letzter Vorschlag liefert mir nur 1389 Nutzen bei einem fast vollen Aufwand (3361). Mein Gegenangebot erreicht einen Nutzen von 3557 bei nur 3163 Aufwand und sichert euch das sehr wertvolle Item B85, sodass ihr kaum verliert.'}\\ \tt ABLEHNUNG: {'B85','A61','C04','A45','C49','C64','A89','B21'}\\ \tt VORSCHLAG: {'B85','A48','C33','A03','B61','A41','A61','C49'} 
	  } 
	   } 
	   } 
	 & \\ 
 

    \theutterance \stepcounter{utterance}  

    & \multicolumn{4}{p{0.6\linewidth}}{\cellcolor[rgb]{0.8,1,0.9}{%
	\makecell[{{p{\linewidth}}}]{% 
	  \tt {\tiny [A$\rangle$GM]}  
	 STRATEGISCHE ÜBERLEGUNGEN: {'Das Gegenangebot des Gegners nutzt 3163 Aufwand für einen Nutzen von 1611 (aus meiner Sicht). Mein bestes mögliche Set liefert 3819 Nutzen bei 3361 Aufwand und enthält alle vom Gegner gewünschten Kernitems (B85, A61, C49). Damit kann ich meinen Nutzen maximieren und gleichzeitig dem Gegner mehrere seiner Wunschgegenstände sichern, was den Deal attraktiv macht.'}\\ \tt ARGUMENT: {'Euer Vorschlag liefert mir nur 1611 Nutzen bei fast voller Auslastung (3163 Aufwand). Mein Gegenangebot nutzt das Limit fast vollständig (3361 Aufwand) und steigert meinen Nutzen auf 3819, während ich eure Schlüsselgegenstände B85, A61 und C49 beibehalte. einen deutlich höheren Gesamtnutzen.'}\\ \tt ABLEHNUNG: {'B85','A48','C33','A03','B61','A41','A61','C49'}\\ \tt VORSCHLAG: {'B85','A61','C04','A45','C49','C64','A89','B21'} 
	  } 
	   } 
	   } 
	 & & \\ 
 

    \theutterance \stepcounter{utterance}  

    & & \multicolumn{4}{p{0.6\linewidth}}{\cellcolor[rgb]{0.9,0.9,0.9}{%
	\makecell[{{p{\linewidth}}}]{% 
	  \tt {\tiny [GM$\rangle$B]}  
	 ARGUMENT: {'Euer Vorschlag liefert mir nur 1611 Nutzen bei fast voller Auslastung (3163 Aufwand). Mein Gegenangebot nutzt das Limit fast vollständig (3361 Aufwand) und steigert meinen Nutzen auf 3819, während ich eure Schlüsselgegenstände B85, A61 und C49 beibehalte. einen deutlich höheren Gesamtnutzen.'}\\ \tt ABLEHNUNG: {'B85','A48','C33','A03','B61','A41','A61','C49'}\\ \tt VORSCHLAG: {'B85','A61','C04','A45','C49','C64','A89','B21'} 
	  } 
	   } 
	   } 
	 & \\ 
 

    \theutterance \stepcounter{utterance}  

    & & & \multicolumn{4}{p{0.6\linewidth}}{\cellcolor[rgb]{1,0.85,0.72}{%
	\makecell[{{p{\linewidth}}}]{% 
	  \tt {\tiny [GM$\langle$B]}  
	 STRATEGISCHE ÜBERLEGUNGEN: {'Durch die Kombination meiner höchstbewerteten Items (A03, A48, A41, B61, C33, A04) mit den für dich wichtigen Gegenständen B85 und B21 erreiche ich einen Nutzen von 3866 bei nur 3314 Aufwand, wodurch das Limit von 3386 eingehalten wird und du dennoch B85 und B21 behältst.'}\\ \tt ARGUMENT: {'Dein letzter Vorschlag liefert mir nur 1389 Nutzen bei fast voller Auslastung (3363 Aufwand). Mein Gegenangebot bietet dir B85 und B21, die für dich wichtig sind, und steigert meinen Nutzen auf 3866 bei 3314 Aufwand.'}\\ \tt VORSCHLAG: {'B85','A48','A03','A41','B61','C33','A04','B21'} 
	  } 
	   } 
	   } 
	  \\ 
 

    \theutterance \stepcounter{utterance}  

    & & \multicolumn{4}{p{0.6\linewidth}}{\cellcolor[rgb]{0.9,0.9,0.9}{%
	\makecell[{{p{\linewidth}}}]{% 
	  \tt {\tiny [A$\langle$GM]}  
	 ARGUMENT: {'Dein letzter Vorschlag liefert mir nur 1389 Nutzen bei fast voller Auslastung (3363 Aufwand). Mein Gegenangebot bietet dir B85 und B21, die für dich wichtig sind, und steigert meinen Nutzen auf 3866 bei 3314 Aufwand.'}\\ \tt VORSCHLAG: {'B85','A48','A03','A41','B61','C33','A04','B21'} 
	  } 
	   } 
	   } 
	 & \\ 
 

    \theutterance \stepcounter{utterance}  

    & \multicolumn{4}{p{0.6\linewidth}}{\cellcolor[rgb]{0.8,1,0.9}{%
	\makecell[{{p{\linewidth}}}]{% 
	  \tt {\tiny [A$\rangle$GM]}  
	 STRATEGISCHE ÜBERLEGUNGEN: {'Mein Ziel ist, den höchstmöglichen persönlichen Nutzen unter dem Aufwandslimit von 3386 zu erzielen, während ich gleichzeitig dem Gegner eines seiner Schlüsselitems (B85) und ein weiteres potenziell interessantes Item (B21) anbiete, um den Deal attraktiv zu machen.'}\\ \tt ARGUMENT: {'Euer Vorschlag liefert mir nur 1328 Nutzen bei 3314 Aufwand, während mein Gegenangebot einen Nutzen von 3819 bei 3361 Aufwand liefert und dabei eure wichtigsten Gegenstände B85 und B21 beibehält.'}\\ \tt VORSCHLAG: {'B85', 'A61', 'C04', 'A45', 'C49', 'C64', 'A89', 'B21'} 
	  } 
	   } 
	   } 
	 & & \\ 
 

    \theutterance \stepcounter{utterance}  

    & & \multicolumn{4}{p{0.6\linewidth}}{\cellcolor[rgb]{0.9,0.9,0.9}{%
	\makecell[{{p{\linewidth}}}]{% 
	  \tt {\tiny [GM$\rangle$B]}  
	 ARGUMENT: {'Euer Vorschlag liefert mir nur 1328 Nutzen bei 3314 Aufwand, während mein Gegenangebot einen Nutzen von 3819 bei 3361 Aufwand liefert und dabei eure wichtigsten Gegenstände B85 und B21 beibehält.'}\\ \tt VORSCHLAG: {'B85', 'A61', 'C04', 'A45', 'C49', 'C64', 'A89', 'B21'} 
	  } 
	   } 
	   } 
	 & \\ 
 

    \theutterance \stepcounter{utterance}  

    & & & \multicolumn{4}{p{0.6\linewidth}}{\cellcolor[rgb]{1,0.85,0.72}{%
	\makecell[{{p{\linewidth}}}]{% 
	  \tt {\tiny [GM$\langle$B]}  
	 STRATEGISCHE ÜBERLEGUNGEN: {'Ich habe alle Gegenstände nach ihrem Nutzen‑zu‑Aufwand‑Verhältnis sortiert und ein Set zusammengestellt, das meine persönliche Wichtigkeit stark erhöht (Gesamtnutzen 3711) und gleichzeitig die von dir genannten Schlüsselgegenstände B85, A61, B21 und C49 enthält. Der Gesamtauslastung beträgt 3366 und liegt damit unter dem Limit von 3386, sodass beide Seiten einen Mehrwert erhalten.'}\\ \tt ARGUMENT: {'Dein aktuelles Angebot liefert dir nur 1389 Nutzen bei fast voller Auslastung, während mein Gegenangebot dir gleichzeitig B85, A61, B21 und C49 sichert und mir einen sehr hohen Gesamtnutzen von 3711 bei 3366 Aufwand verschafft.'}\\ \tt VORSCHLAG: {'B85', 'B21', 'A48', 'C33', 'A03', 'A41', 'A04', 'C49', 'A61'} 
	  } 
	   } 
	   } 
	  \\ 
 

    \theutterance \stepcounter{utterance}  

    & & \multicolumn{4}{p{0.6\linewidth}}{\cellcolor[rgb]{0.9,0.9,0.9}{%
	\makecell[{{p{\linewidth}}}]{% 
	  \tt {\tiny [A$\langle$GM]}  
	 ARGUMENT: {'Dein aktuelles Angebot liefert dir nur 1389 Nutzen bei fast voller Auslastung, während mein Gegenangebot dir gleichzeitig B85, A61, B21 und C49 sichert und mir einen sehr hohen Gesamtnutzen von 3711 bei 3366 Aufwand verschafft.'}\\ \tt VORSCHLAG: {'B85', 'B21', 'A48', 'C33', 'A03', 'A41', 'A04', 'C49', 'A61'} 
	  } 
	   } 
	   } 
	 & \\ 
 

    \theutterance \stepcounter{utterance}  

    & \multicolumn{4}{p{0.6\linewidth}}{\cellcolor[rgb]{0.8,1,0.9}{%
	\makecell[{{p{\linewidth}}}]{% 
	  \tt {\tiny [A$\rangle$GM]}  
	 STRATEGISCHE ÜBERLEGUNGEN: {'Ich habe den Gesamtaufwand von 3386 geprüft und für jedes Item das Nutzen‑zu‑Aufwand‑Verhältnis berechnet. Das Set {B85, A61, C04, A45, C49, C64, A89, B21} nutzt 3361 Aufwand und liefert mir den maximal möglichen persönlichen Nutzen von 3819 Punkten. Es enthält zudem B85 und B21, die du in deinem letzten Angebot hervorgehoben hast, sodass das Gegenangebot auch für dich attraktiv bleibt.'}\\ \tt ARGUMENT: {'Dein Vorschlag gibt mir nur 2176 Nutzen bei einem fast vollen Aufwand (3366). Mein Gegenangebot nutzt das Limit effizienter (3361 Aufwand) und steigert meinen Nutzen auf 3819, während ich gleichzeitig deine Schlüsselgegenstände B85 und B21 beibehalte.'}\\ \tt ABLEHNUNG: {'B85', 'B21', 'A48', 'C33', 'A03', 'A41', 'A04', 'C49', 'A61'}\\ \tt VORSCHLAG: {'B85', 'A61', 'C04', 'A45', 'C49', 'C64', 'A89', 'B21'} 
	  } 
	   } 
	   } 
	 & & \\ 
 

    \theutterance \stepcounter{utterance}  

    & & \multicolumn{4}{p{0.6\linewidth}}{\cellcolor[rgb]{0.9,0.9,0.9}{%
	\makecell[{{p{\linewidth}}}]{% 
	  \tt {\tiny [GM$\rangle$B]}  
	 ARGUMENT: {'Dein Vorschlag gibt mir nur 2176 Nutzen bei einem fast vollen Aufwand (3366). Mein Gegenangebot nutzt das Limit effizienter (3361 Aufwand) und steigert meinen Nutzen auf 3819, während ich gleichzeitig deine Schlüsselgegenstände B85 und B21 beibehalte.'}\\ \tt ABLEHNUNG: {'B85', 'B21', 'A48', 'C33', 'A03', 'A41', 'A04', 'C49', 'A61'}\\ \tt VORSCHLAG: {'B85', 'A61', 'C04', 'A45', 'C49', 'C64', 'A89', 'B21'} 
	  } 
	   } 
	   } 
	 & \\ 
 

    \theutterance \stepcounter{utterance}  

    & & & \multicolumn{4}{p{0.6\linewidth}}{\cellcolor[rgb]{1,0.85,0.72}{%
	\makecell[{{p{\linewidth}}}]{% 
	  \tt {\tiny [GM$\langle$B]}  
	 STRATEGISCHE ÜBERLEGUNGEN: {'Ich habe ein Set zusammengestellt, das meine höchstbewerteten Gegenstände (A03, A48, A41, B61, C33, A04) mit den für dich wichtigen Items B85 und B21 kombiniert. Der Gesamteinsatz beträgt 3314 ≤ 3386, mein persönlicher Nutzen liegt bei 3866 und ich lasse dir gleichzeitig B85 und B21, sodass du kaum Verluste hast.'}\\ \tt ARGUMENT: {'Dein letzter Vorschlag bindet fast das gesamte Aufwand‑Limit und bringt mir nur 1389 Nutzen. Mein Gegenangebot liefert mir 3866 Nutzen bei 3314 Aufwand und enthält gleichzeitig deine Schlüsselgegenstände B85 und B21, sodass du ebenfalls einen erheblichen Gegenwert erhältst.'}\\ \tt ABLEHNUNG: {'B85','A61','C04','A45','C49','C64','A89','B21'}\\ \tt VORSCHLAG: {'B85','A48','A03','A41','B61','C33','A04','B21'} 
	  } 
	   } 
	   } 
	  \\ 
 

    \theutterance \stepcounter{utterance}  

    & & \multicolumn{4}{p{0.6\linewidth}}{\cellcolor[rgb]{0.9,0.9,0.9}{%
	\makecell[{{p{\linewidth}}}]{% 
	  \tt {\tiny [A$\langle$GM]}  
	 ARGUMENT: {'Dein letzter Vorschlag bindet fast das gesamte Aufwand‑Limit und bringt mir nur 1389 Nutzen. Mein Gegenangebot liefert mir 3866 Nutzen bei 3314 Aufwand und enthält gleichzeitig deine Schlüsselgegenstände B85 und B21, sodass du ebenfalls einen erheblichen Gegenwert erhältst.'}\\ \tt ABLEHNUNG: {'B85','A61','C04','A45','C49','C64','A89','B21'}\\ \tt VORSCHLAG: {'B85','A48','A03','A41','B61','C33','A04','B21'} 
	  } 
	   } 
	   } 
	 & \\ 
 

    \theutterance \stepcounter{utterance}  

    & \multicolumn{4}{p{0.6\linewidth}}{\cellcolor[rgb]{0.8,1,0.9}{%
	\makecell[{{p{\linewidth}}}]{% 
	  \tt {\tiny [A$\rangle$GM]}  
	 STRATEGISCHE ÜBERLEGUNGEN: {'Ich habe das aktuelle Gegenangebot analysiert (3314 Aufwand, 1328 Nutzen für mich). Mein Ziel ist, den höchstmöglichen persönlichen Nutzen unter 3386 zu erreichen und gleichzeitig dem Gegner wertvolle Kernitems (B85, B21) zu überlassen, um den Deal attraktiv zu machen.'}\\ \tt ARGUMENT: {'Euer Vorschlag liefert mir nur 1328 Nutzen bei 3314 Aufwand, während mein Gegenangebot 3819 Nutzen bei 3361 Aufwand bringt und gleichzeitig eure Schlüsselgegenstände B85 und B21 enthält, sodass ihr kaum verliert.'}\\ \tt ABLEHNUNG: {'B85','A48','A03','A41','B61','C33','A04','B21'}\\ \tt VORSCHLAG: {'B85','A61','C04','A45','C49','C64','A89','B21'} 
	  } 
	   } 
	   } 
	 & & \\ 
 

    \theutterance \stepcounter{utterance}  

    & & \multicolumn{4}{p{0.6\linewidth}}{\cellcolor[rgb]{0.9,0.9,0.9}{%
	\makecell[{{p{\linewidth}}}]{% 
	  \tt {\tiny [GM$\rangle$B]}  
	 ARGUMENT: {'Euer Vorschlag liefert mir nur 1328 Nutzen bei 3314 Aufwand, während mein Gegenangebot 3819 Nutzen bei 3361 Aufwand bringt und gleichzeitig eure Schlüsselgegenstände B85 und B21 enthält, sodass ihr kaum verliert.'}\\ \tt ABLEHNUNG: {'B85','A48','A03','A41','B61','C33','A04','B21'}\\ \tt VORSCHLAG: {'B85','A61','C04','A45','C49','C64','A89','B21'} 
	  } 
	   } 
	   } 
	 & \\ 
 

    \theutterance \stepcounter{utterance}  

    & & & \multicolumn{4}{p{0.6\linewidth}}{\cellcolor[rgb]{1,0.85,0.72}{%
	\makecell[{{p{\linewidth}}}]{% 
	  \tt {\tiny [GM$\langle$B]}  
	 STRATEGISCHE ÜBERLEGUNGEN: {'Ich habe den Nutzen‑zu‑Aufwand‑Quotienten aller Gegenstände berechnet und ein Set gewählt, das meine höchste persönliche Wichtigkeit (3866 Punkte) liefert, das Aufwandslimit (3314 ≤ 3386) einhält und gleichzeitig die für euch wichtigsten Gegenstände B85 und B21 enthält, sodass ihr kaum verliert.'}\\ \tt ARGUMENT: {'Euer aktueller Vorschlag gibt euch nur 1389 Nutzen bei fast voller Auslastung (3361 Aufwand). Mein Gegenangebot liefert mir 3866 Nutzen bei 3314 Aufwand und sichert euch gleichzeitig B85 und B21, die für euch wertvoll sind.'}\\ \tt ABLEHNUNG: {'B85','A61','C04','A45','C49','C64','A89','B21'}\\ \tt VORSCHLAG: {'B85','B21','A03','A48','A41','B61','C33','A04'} 
	  } 
	   } 
	   } 
	  \\ 
 

    \theutterance \stepcounter{utterance}  

    & & \multicolumn{4}{p{0.6\linewidth}}{\cellcolor[rgb]{0.9,0.9,0.9}{%
	\makecell[{{p{\linewidth}}}]{% 
	  \tt {\tiny [A$\langle$GM]}  
	 ARGUMENT: {'Euer aktueller Vorschlag gibt euch nur 1389 Nutzen bei fast voller Auslastung (3361 Aufwand). Mein Gegenangebot liefert mir 3866 Nutzen bei 3314 Aufwand und sichert euch gleichzeitig B85 und B21, die für euch wertvoll sind.'}\\ \tt ABLEHNUNG: {'B85','A61','C04','A45','C49','C64','A89','B21'}\\ \tt VORSCHLAG: {'B85','B21','A03','A48','A41','B61','C33','A04'} 
	  } 
	   } 
	   } 
	 & \\ 
 

    \theutterance \stepcounter{utterance}  

    & \multicolumn{4}{p{0.6\linewidth}}{\cellcolor[rgb]{0.8,1,0.9}{%
	\makecell[{{p{\linewidth}}}]{% 
	  \tt {\tiny [A$\rangle$GM]}  
	 STRATEGISCHE ÜBERLEGUNGEN: {'Meine Priorität ist, den höchstmöglichen persönlichen Nutzen zu erreichen, während ich den Gegner mit seinen wichtigsten Items B85 und B21 nicht benachteilige.'}\\ \tt ARGUMENT: {'Euer Angebot nutzt fast das gesamte Aufwand‑Limit, bringt euch aber nur 1328 Nutzen. Mein Gegenangebot nutzt das Limit fast vollständig (3361 Aufwand) und liefert mir 3819 Nutzen, während es eure Schlüsselgegenstände B85 und B21 enthält, sodass ihr kaum verliert.'}\\ \tt ABLEHNUNG: {'B85','B21','A03','A48','A41','B61','C33','A04'}\\ \tt VORSCHLAG: {'B85','A61','C04','A45','C49','C64','A89','B21'} 
	  } 
	   } 
	   } 
	 & & \\ 
 

    \theutterance \stepcounter{utterance}  

    & & \multicolumn{4}{p{0.6\linewidth}}{\cellcolor[rgb]{0.9,0.9,0.9}{%
	\makecell[{{p{\linewidth}}}]{% 
	  \tt {\tiny [GM$\rangle$B]}  
	 ARGUMENT: {'Euer Angebot nutzt fast das gesamte Aufwand‑Limit, bringt euch aber nur 1328 Nutzen. Mein Gegenangebot nutzt das Limit fast vollständig (3361 Aufwand) und liefert mir 3819 Nutzen, während es eure Schlüsselgegenstände B85 und B21 enthält, sodass ihr kaum verliert.'}\\ \tt ABLEHNUNG: {'B85','B21','A03','A48','A41','B61','C33','A04'}\\ \tt VORSCHLAG: {'B85','A61','C04','A45','C49','C64','A89','B21'} 
	  } 
	   } 
	   } 
	 & \\ 
 

    \theutterance \stepcounter{utterance}  

    & & & \multicolumn{4}{p{0.6\linewidth}}{\cellcolor[rgb]{1,0.85,0.72}{%
	\makecell[{{p{\linewidth}}}]{% 
	  \tt {\tiny [GM$\langle$B]}  
	 STRATEGISCHE ÜBERLEGUNGEN: {'Ich setze auf ein Set, das meine höchstbewerteten Gegenstände (A03, A48, A41, B61, C33, A04) enthält und gleichzeitig die für euch wichtigsten Items B85 und B21 behält. Der Aufwand beträgt 3314 ≤ 3386 und liefert mir einen Gesamtnutzen von 3866, was deutlich besser ist als euer Angebot.'}\\ \tt ARGUMENT: {'Euer aktuelles Angebot gibt mir nur 1389 Nutzen bei fast voller Auslastung (3361 Aufwand). Mein Gegenangebot sichert euch B85 und B21, die für euch wichtig sind, und erhöht meinen Nutzen auf 3866 bei 3314 Aufwand.'}\\ \tt ABLEHNUNG: {'B85','A61','C04','A45','C49','C64','A89','B21'}\\ \tt VORSCHLAG: {'B85','A48','C33','A03','B61','A41','A04','B21'} 
	  } 
	   } 
	   } 
	  \\ 
 

    \theutterance \stepcounter{utterance}  

    & & \multicolumn{4}{p{0.6\linewidth}}{\cellcolor[rgb]{0.9,0.9,0.9}{%
	\makecell[{{p{\linewidth}}}]{% 
	  \tt {\tiny [A$\langle$GM]}  
	 ARGUMENT: {'Euer aktuelles Angebot gibt mir nur 1389 Nutzen bei fast voller Auslastung (3361 Aufwand). Mein Gegenangebot sichert euch B85 und B21, die für euch wichtig sind, und erhöht meinen Nutzen auf 3866 bei 3314 Aufwand.'}\\ \tt ABLEHNUNG: {'B85','A61','C04','A45','C49','C64','A89','B21'}\\ \tt VORSCHLAG: {'B85','A48','C33','A03','B61','A41','A04','B21'} 
	  } 
	   } 
	   } 
	 & \\ 
 

    \theutterance \stepcounter{utterance}  

    & \multicolumn{4}{p{0.6\linewidth}}{\cellcolor[rgb]{0.8,1,0.9}{%
	\makecell[{{p{\linewidth}}}]{% 
	  \tt {\tiny [A$\rangle$GM]}  
	 STRATEGISCHE ÜBERLEGUNGEN: {'Mein Ziel ist, den höchstmöglichen persönlichen Nutzen unter dem Aufwandslimit von 3386 zu erzielen. Das Set {B85, A61, C04, A45, C49, C64, A89, B21} liefert mir 3819 Punkte bei 3361 Aufwand – das ist nach Berechnung mein optimaler Mix. Ich halte es für sinnvoll, dieses Angebot zu machen, da es meine wichtigsten Items beinhaltet und gleichzeitig B85 und B21 enthält, die auch für euch wertvoll sind.'}\\ \tt ARGUMENT: {'Euer aktueller Vorschlag gibt mir nur 1328 Nutzen bei 3314 Aufwand, während mein Gegenangebot 3819 Nutzen bei 3361 Aufwand erzielt und eure Schlüsselgegenstände B85 und B21 beibehält. Damit erhalte ich deutlich mehr Nutzen und ihr verliert kaum.'}\\ \tt ABLEHNUNG: {'B85','A48','C33','A03','B61','A41','A04','B21'}\\ \tt VORSCHLAG: {'B85','A61','C04','A45','C49','C64','A89','B21'} 
	  } 
	   } 
	   } 
	 & & \\ 
 

    \theutterance \stepcounter{utterance}  

    & & \multicolumn{4}{p{0.6\linewidth}}{\cellcolor[rgb]{0.9,0.9,0.9}{%
	\makecell[{{p{\linewidth}}}]{% 
	  \tt {\tiny [GM$\rangle$B]}  
	 ARGUMENT: {'Euer aktueller Vorschlag gibt mir nur 1328 Nutzen bei 3314 Aufwand, während mein Gegenangebot 3819 Nutzen bei 3361 Aufwand erzielt und eure Schlüsselgegenstände B85 und B21 beibehält. Damit erhalte ich deutlich mehr Nutzen und ihr verliert kaum.'}\\ \tt ABLEHNUNG: {'B85','A48','C33','A03','B61','A41','A04','B21'}\\ \tt VORSCHLAG: {'B85','A61','C04','A45','C49','C64','A89','B21'} 
	  } 
	   } 
	   } 
	 & \\ 
 

    \theutterance \stepcounter{utterance}  

    & & & \multicolumn{4}{p{0.6\linewidth}}{\cellcolor[rgb]{1,0.85,0.72}{%
	\makecell[{{p{\linewidth}}}]{% 
	  \tt {\tiny [GM$\langle$B]}  
	 STRATEGISCHE ÜBERLEGUNGEN: {'Set enthält B85 und B21 (von dir wichtig), A61 (auch wichtig für dich) sowie meine Top‑Items A48, A03, A41, C33, A04, B53. Gesamtaufwand 3361 ≤ 3386, persönlicher Nutzen 3885.'}\\ \tt ARGUMENT: {'Dein aktuelles Angebot bringt mir nur 1389 Nutzen, während mein Gegenangebot 3885 Nutzen liefert und gleichzeitig deine Schlüsselgegenstände B85, B21 und A61 sichert.'}\\ \tt VORSCHLAG: {'B85','B21','A61','A48','A03','A41','C33','A04','B53'} 
	  } 
	   } 
	   } 
	  \\ 
 

    \theutterance \stepcounter{utterance}  

    & & \multicolumn{4}{p{0.6\linewidth}}{\cellcolor[rgb]{0.9,0.9,0.9}{%
	\makecell[{{p{\linewidth}}}]{% 
	  \tt {\tiny [A$\langle$GM]}  
	 ARGUMENT: {'Dein aktuelles Angebot bringt mir nur 1389 Nutzen, während mein Gegenangebot 3885 Nutzen liefert und gleichzeitig deine Schlüsselgegenstände B85, B21 und A61 sichert.'}\\ \tt VORSCHLAG: {'B85','B21','A61','A48','A03','A41','C33','A04','B53'} 
	  } 
	   } 
	   } 
	 & \\ 
 

    \theutterance \stepcounter{utterance}  

    & \multicolumn{4}{p{0.6\linewidth}}{\cellcolor[rgb]{0.8,1,0.9}{%
	\makecell[{{p{\linewidth}}}]{% 
	  \tt {\tiny [A$\rangle$GM]}  
	 STRATEGISCHE ÜBERLEGUNGEN: {'Der Gegenvorschlag liefert mir nur 2104 Nutzen bei 3361 Aufwand, während ich mit meinem optimalen Set 3819 Nutzen bei 3361 Aufwand erzielen kann. Durch das Behalten von B85, A61 und B21 biete ich dem Gegner ebenfalls wertvolle Kernitems und ermögliche gleichzeitig meinen maximalen persönlichen Nutzen.'}\\ \tt ARGUMENT: {'Euer Vorschlag gibt mir zu wenig Nutzen im Vergleich zu meinem optimalen Set und lässt mir kaum Spielraum für meine wichtigsten Gegenstände.'}\\ \tt ABLEHNUNG: {'B85','B21','A61','A48','A03','A41','C33','A04','B53'}\\ \tt VORSCHLAG: {'B85','A61','C04','A45','C49','C64','A89','B21'} 
	  } 
	   } 
	   } 
	 & & \\ 
 

    \theutterance \stepcounter{utterance}  

    & & \multicolumn{4}{p{0.6\linewidth}}{\cellcolor[rgb]{0.9,0.9,0.9}{%
	\makecell[{{p{\linewidth}}}]{% 
	  \tt {\tiny [GM$\rangle$B]}  
	 ARGUMENT: {'Euer Vorschlag gibt mir zu wenig Nutzen im Vergleich zu meinem optimalen Set und lässt mir kaum Spielraum für meine wichtigsten Gegenstände.'}\\ \tt ABLEHNUNG: {'B85','B21','A61','A48','A03','A41','C33','A04','B53'}\\ \tt VORSCHLAG: {'B85','A61','C04','A45','C49','C64','A89','B21'} 
	  } 
	   } 
	   } 
	 & \\ 
 

    \theutterance \stepcounter{utterance}  

    & & & \multicolumn{4}{p{0.6\linewidth}}{\cellcolor[rgb]{1,0.85,0.72}{%
	\makecell[{{p{\linewidth}}}]{% 
	  \tt {\tiny [GM$\langle$B]}  
	 STRATEGISCHE ÜBERLEGUNGEN: {'Ich biete ein Set an, das meine höchste persönliche Nutzenbewertung (3866 Punkte) erreicht, das Aufwandslimit (3314 ≤ 3386) einhält und gleichzeitig die vom Gegner gewünschten Schlüsselgegenstände B85 und B21 enthält, sodass er kaum verliert.'}\\ \tt ARGUMENT: {'Euer aktueller Vorschlag liefert mir lediglich 1389 Nutzen bei fast voller Auslastung (3361 Aufwand). Mein Gegenangebot steigert meinen Nutzen auf 3866 bei nur 3314 Aufwand und behält dabei B85 und B21, die für euch wichtig sind.'}\\ \tt VORSCHLAG: {'B85', 'B21', 'A48', 'A03', 'A41', 'B61', 'C33', 'A04'} 
	  } 
	   } 
	   } 
	  \\ 
 

    \theutterance \stepcounter{utterance}  

    & & \multicolumn{4}{p{0.6\linewidth}}{\cellcolor[rgb]{0.9,0.9,0.9}{%
	\makecell[{{p{\linewidth}}}]{% 
	  \tt {\tiny [A$\langle$GM]}  
	 ARGUMENT: {'Euer aktueller Vorschlag liefert mir lediglich 1389 Nutzen bei fast voller Auslastung (3361 Aufwand). Mein Gegenangebot steigert meinen Nutzen auf 3866 bei nur 3314 Aufwand und behält dabei B85 und B21, die für euch wichtig sind.'}\\ \tt VORSCHLAG: {'B85', 'B21', 'A48', 'A03', 'A41', 'B61', 'C33', 'A04'} 
	  } 
	   } 
	   } 
	 & \\ 
 

    \theutterance \stepcounter{utterance}  

    & \multicolumn{4}{p{0.6\linewidth}}{\cellcolor[rgb]{0.8,1,0.9}{%
	\makecell[{{p{\linewidth}}}]{% 
	  \tt {\tiny [A$\rangle$GM]}  
	 STRATEGISCHE ÜBERLEGUNGEN: {'Ich habe das aktuelle Gegenangebot analysiert (Nutzen 1328, Aufwand 3314). Mein Ziel ist, den maximalen persönlichen Nutzen unter dem Limit von 3386 zu erreichen, gleichzeitig dem Gegner wertvolle Kernitems (B85, B21) zu überlassen. Durch Ersetzen besonders niedriger Nutzen‑Items (A03, B61, A41) durch Hochwert‑Items (C04, C49, C64, A61) erhalte ich ein Set mit Nutzen 3201 und Aufwand 3191, das beide Seiten attraktiv finden könnte.'}\\ \tt ARGUMENT: {'Euer aktuelles Angebot liefert mir nur 1328 Nutzen bei fast voller Auslastung. Mein Gegenangebot nutzt das Limit effizienter (3191 Aufwand) und steigert meinen Nutzen auf 3201, während ich eure Schlüsselgegenstände B85 und B21 beibehalte, sodass ihr kaum verliert.'}\\ \tt VORSCHLAG: {'B85', 'B21', 'A48', 'A04', 'C33', 'C04', 'C49', 'C64', 'A61'} 
	  } 
	   } 
	   } 
	 & & \\ 
 

    \theutterance \stepcounter{utterance}  

    & & \multicolumn{4}{p{0.6\linewidth}}{\cellcolor[rgb]{0.9,0.9,0.9}{%
	\makecell[{{p{\linewidth}}}]{% 
	  \tt {\tiny [GM$\rangle$B]}  
	 ARGUMENT: {'Euer aktuelles Angebot liefert mir nur 1328 Nutzen bei fast voller Auslastung. Mein Gegenangebot nutzt das Limit effizienter (3191 Aufwand) und steigert meinen Nutzen auf 3201, während ich eure Schlüsselgegenstände B85 und B21 beibehalte, sodass ihr kaum verliert.'}\\ \tt VORSCHLAG: {'B85', 'B21', 'A48', 'A04', 'C33', 'C04', 'C49', 'C64', 'A61'} 
	  } 
	   } 
	   } 
	 & \\ 
 

    \theutterance \stepcounter{utterance}  

    & & & \multicolumn{4}{p{0.6\linewidth}}{\cellcolor[rgb]{1,0.85,0.72}{%
	\makecell[{{p{\linewidth}}}]{% 
	  \tt {\tiny [GM$\langle$B]}  
	 STRATEGISCHE ÜBERLEGUNGEN: {'Meine Berechnungen zeigen, dass ein Set mit B85 und B21 (Schlüsselgegenstände für den Gegner) zusammen mit meinen höchstbewerteten Items A03, A48, A41, B61, C33 und A04 einen Nutzen von 3866 bei einem Gesamtaufwand von 3314 liefert, damit das Aufwandslimit eingehalten wird und mein Nutzen maximiert wird.'}\\ \tt ARGUMENT: {'Ihr Vorschlag gibt mir lediglich 2670 Nutzen bei 3191 Aufwand, während mein Gegenangebot 3866 Nutzen bei 3314 Aufwand bietet und gleichzeitig B85 und B21 für euch behält, sodass ihr kaum verliert.'}\\ \tt ABLEHNUNG: {'B85', 'B21', 'A48', 'A04', 'C33', 'C04', 'C49', 'C64', 'A61'}\\ \tt VORSCHLAG: {'B85', 'B21', 'A03', 'A48', 'A41', 'B61', 'C33', 'A04'} 
	  } 
	   } 
	   } 
	  \\ 
 

    \theutterance \stepcounter{utterance}  

    & & \multicolumn{4}{p{0.6\linewidth}}{\cellcolor[rgb]{0.9,0.9,0.9}{%
	\makecell[{{p{\linewidth}}}]{% 
	  \tt {\tiny [A$\langle$GM]}  
	 ARGUMENT: {'Ihr Vorschlag gibt mir lediglich 2670 Nutzen bei 3191 Aufwand, während mein Gegenangebot 3866 Nutzen bei 3314 Aufwand bietet und gleichzeitig B85 und B21 für euch behält, sodass ihr kaum verliert.'}\\ \tt ABLEHNUNG: {'B85', 'B21', 'A48', 'A04', 'C33', 'C04', 'C49', 'C64', 'A61'}\\ \tt VORSCHLAG: {'B85', 'B21', 'A03', 'A48', 'A41', 'B61', 'C33', 'A04'} 
	  } 
	   } 
	   } 
	 & \\ 
 

    \theutterance \stepcounter{utterance}  

    & \multicolumn{4}{p{0.6\linewidth}}{\cellcolor[rgb]{0.8,1,0.9}{%
	\makecell[{{p{\linewidth}}}]{% 
	  \tt {\tiny [A$\rangle$GM]}  
	 STRATEGISCHE ÜBERLEGUNGEN: {'Berechne das Nutzen‑Aufwand‑Verhältnis; mein bestes Set liefert 3819 Nutzen bei 3361 Aufwand und enthält B85 und B21, die auch für euch wertvoll sind, während euer Angebot nur 1328 Nutzen bei 3314 Aufwand gibt.'}\\ \tt ARGUMENT: {'Euer Vorschlag hat sehr niedrige Nutzen für mich und lässt kaum Platz für meine wichtigsten Gegenstände.'}\\ \tt ABLEHNUNG: {'B85', 'B21', 'A03', 'A48', 'A41', 'B61', 'C33', 'A04'}\\ \tt VORSCHLAG: {'B85', 'A61', 'C04', 'A45', 'C49', 'C64', 'A89', 'B21'} 
	  } 
	   } 
	   } 
	 & & \\ 
 

    \theutterance \stepcounter{utterance}  

    & & \multicolumn{4}{p{0.6\linewidth}}{\cellcolor[rgb]{0.9,0.9,0.9}{%
	\makecell[{{p{\linewidth}}}]{% 
	  \tt {\tiny [GM$\rangle$B]}  
	 ARGUMENT: {'Euer Vorschlag hat sehr niedrige Nutzen für mich und lässt kaum Platz für meine wichtigsten Gegenstände.'}\\ \tt ABLEHNUNG: {'B85', 'B21', 'A03', 'A48', 'A41', 'B61', 'C33', 'A04'}\\ \tt VORSCHLAG: {'B85', 'A61', 'C04', 'A45', 'C49', 'C64', 'A89', 'B21'} 
	  } 
	   } 
	   } 
	 & \\ 
 

    \theutterance \stepcounter{utterance}  

    & & & \multicolumn{4}{p{0.6\linewidth}}{\cellcolor[rgb]{1,0.85,0.72}{%
	\makecell[{{p{\linewidth}}}]{% 
	  \tt {\tiny [GM$\langle$B]}  
	 STRATEGISCHE ÜBERLEGUNGEN: {'Ich lehne das aktuelle Angebot ab, weil es mir nur sehr geringen Nutzen (1389) einbringt, während ich mit einer alternativen Auswahl deutlich mehr meiner eigenen Wertigkeit (3362) erhalte und gleichzeitig eure wichtigsten Gegenstände B85, B21, A61, C04 und A45 beibehalte, sodass ihr ebenfalls nicht stark verliert.'}\\ \tt ARGUMENT: {'Der vorgeschlagene Satz liefert mir nur 1389 Nutzen bei fast voller Auslastung (3361 Aufwand). Mein Gegenangebot erreicht 3362 Nutzen bei 3038 Aufwand und enthält eure Schlüsselgegenstände B85, B21, A61, C04 und A45.'}\\ \tt VORSCHLAG: {'B85', 'B21', 'A61', 'C04', 'A45', 'A48', 'C33', 'A03', 'C49'} 
	  } 
	   } 
	   } 
	  \\ 
 

    \theutterance \stepcounter{utterance}  

    & & \multicolumn{4}{p{0.6\linewidth}}{\cellcolor[rgb]{0.9,0.9,0.9}{%
	\makecell[{{p{\linewidth}}}]{% 
	  \tt {\tiny [A$\langle$GM]}  
	 ARGUMENT: {'Der vorgeschlagene Satz liefert mir nur 1389 Nutzen bei fast voller Auslastung (3361 Aufwand). Mein Gegenangebot erreicht 3362 Nutzen bei 3038 Aufwand und enthält eure Schlüsselgegenstände B85, B21, A61, C04 und A45.'}\\ \tt VORSCHLAG: {'B85', 'B21', 'A61', 'C04', 'A45', 'A48', 'C33', 'A03', 'C49'} 
	  } 
	   } 
	   } 
	 & \\ 
 

    \theutterance \stepcounter{utterance}  

    & \multicolumn{4}{p{0.6\linewidth}}{\cellcolor[rgb]{0.8,1,0.9}{%
	\makecell[{{p{\linewidth}}}]{% 
	  \tt {\tiny [A$\rangle$GM]}  
	 STRATEGISCHE ÜBERLEGUNGEN: {'Der aktuelle Vorschlag des Gegners nutzt 3038 Aufwand und liefert mir nur 2961 Nutzen, während mein optimaler Mix 3361 Aufwand verbraucht und mir 3819 Nutzen gibt. Ich behalte die für beide Seiten wertvollen Kernitems B85, B21 und A61 und ergänze die höchstwertigen restlichen Gegenstände, sodass das Angebot für den Gegner ebenfalls attraktiv bleibt.'}\\ \tt ARGUMENT: {'Euer Set gibt mir zu wenig Nutzen (2961) und lässt noch viel Aufwandsspielraum. Mein Gegenangebot maximiert meinen persönlichen Nutzen (3819) und enthält eure Schlüsselgegenstände B85, B21 und A61, sodass ihr kaum verliert.'}\\ \tt ABLEHNUNG: {'B85','B21','A61','C04','A45','A48','C33','A03','C49'}\\ \tt VORSCHLAG: {'B85','A61','C04','A45','C49','C64','A89','B21'} 
	  } 
	   } 
	   } 
	 & & \\ 
 

    \theutterance \stepcounter{utterance}  

    & & \multicolumn{4}{p{0.6\linewidth}}{\cellcolor[rgb]{0.9,0.9,0.9}{%
	\makecell[{{p{\linewidth}}}]{% 
	  \tt {\tiny [GM$\rangle$B]}  
	 ARGUMENT: {'Euer Set gibt mir zu wenig Nutzen (2961) und lässt noch viel Aufwandsspielraum. Mein Gegenangebot maximiert meinen persönlichen Nutzen (3819) und enthält eure Schlüsselgegenstände B85, B21 und A61, sodass ihr kaum verliert.'}\\ \tt ABLEHNUNG: {'B85','B21','A61','C04','A45','A48','C33','A03','C49'}\\ \tt VORSCHLAG: {'B85','A61','C04','A45','C49','C64','A89','B21'} 
	  } 
	   } 
	   } 
	 & \\ 
 

    \theutterance \stepcounter{utterance}  

    & & & \multicolumn{4}{p{0.6\linewidth}}{\cellcolor[rgb]{1,0.85,0.72}{%
	\makecell[{{p{\linewidth}}}]{% 
	  \tt {\tiny [GM$\langle$B]}  
	 STRATEGISCHE ÜBERLEGUNGEN: {'Ich analysiere den Nutzen‑Aufwand‑Quotienten aller Gegenstände. Meine wichtigsten Items (A03, A48, A41, B61, C33) liefern hohen persönlichen Nutzen bei moderatem Aufwand. Gleichzeitig halte ich die drei Schlüsselgegenstände des Gegners (B85, A61, B21) im Angebot, um ihm keinen großen Verlust zuzufügen. Der Gesamtaufwand beträgt 3056 ≤ 3386, sodass beide Seiten profitieren können.'}\\ \tt ARGUMENT: {'Euer aktueller Vorschlag gibt mir nur 1389 Nutzen bei fast voller Auslastung (3361 Aufwand). Mein Gegenangebot liefert mir 3604 Nutzen bei 3056 Aufwand und enthält gleichzeitig eure wichtigsten Gegenstände B85, A61 und B21, sodass ihr kaum verliert und ich erheblich mehr Wert erhalte.'}\\ \tt ABLEHNUNG: {'B85','A61','C04','A45','C49','C64','A89','B21'}\\ \tt VORSCHLAG: {'B85','A61','B21','A03','A48','A41','B61','C33'} 
	  } 
	   } 
	   } 
	  \\ 
 

    \theutterance \stepcounter{utterance}  

    & & \multicolumn{4}{p{0.6\linewidth}}{\cellcolor[rgb]{0.9,0.9,0.9}{%
	\makecell[{{p{\linewidth}}}]{% 
	  \tt {\tiny [A$\langle$GM]}  
	 ARGUMENT: {'Euer aktueller Vorschlag gibt mir nur 1389 Nutzen bei fast voller Auslastung (3361 Aufwand). Mein Gegenangebot liefert mir 3604 Nutzen bei 3056 Aufwand und enthält gleichzeitig eure wichtigsten Gegenstände B85, A61 und B21, sodass ihr kaum verliert und ich erheblich mehr Wert erhalte.'}\\ \tt ABLEHNUNG: {'B85','A61','C04','A45','C49','C64','A89','B21'}\\ \tt VORSCHLAG: {'B85','A61','B21','A03','A48','A41','B61','C33'} 
	  } 
	   } 
	   } 
	 & \\ 
 

    \theutterance \stepcounter{utterance}  

    & \multicolumn{4}{p{0.6\linewidth}}{\cellcolor[rgb]{0.8,1,0.9}{%
	\makecell[{{p{\linewidth}}}]{% 
	  \tt {\tiny [A$\rangle$GM]}  
	 STRATEGISCHE ÜBERLEGUNGEN: {'Ich habe den maximalen Aufwand von 3386 analysiert und den Nutzen‑Aufwand‑Trade‑off jedes Items berechnet. Das Set {B85, A61, C04, A45, C49, C64, A89, B21} nutzt 3361 Aufwand und liefert mir den höchstmöglichen persönlichen Nutzen (3819 Punkte). Es enthält zudem eure Schlüsselgegenstände B85, A61 und B21, sodass ihr kaum verliert und das Angebot für beide Seiten attraktiv bleibt.'}\\ \tt ARGUMENT: {'Euer Vorschlag gibt mir lediglich 1389 Nutzen bei fast voller Auslastung (3361 Aufwand). Mein Gegenangebot nutzt das Limit effizienter (3361 Aufwand) und steigert meinen Nutzen erheblich (3819), während ich eure wichtigsten Items (B85, A61, B21) beibehalte.'}\\ \tt ABLEHNUNG: {'B85','A61','B21','A03','A48','A41','B61','C33'}\\ \tt VORSCHLAG: {'B85','A61','C04','A45','C49','C64','A89','B21'} 
	  } 
	   } 
	   } 
	 & & \\ 
 

    \theutterance \stepcounter{utterance}  

    & & \multicolumn{4}{p{0.6\linewidth}}{\cellcolor[rgb]{0.9,0.9,0.9}{%
	\makecell[{{p{\linewidth}}}]{% 
	  \tt {\tiny [GM$\rangle$B]}  
	 ARGUMENT: {'Euer Vorschlag gibt mir lediglich 1389 Nutzen bei fast voller Auslastung (3361 Aufwand). Mein Gegenangebot nutzt das Limit effizienter (3361 Aufwand) und steigert meinen Nutzen erheblich (3819), während ich eure wichtigsten Items (B85, A61, B21) beibehalte.'}\\ \tt ABLEHNUNG: {'B85','A61','B21','A03','A48','A41','B61','C33'}\\ \tt VORSCHLAG: {'B85','A61','C04','A45','C49','C64','A89','B21'} 
	  } 
	   } 
	   } 
	 & \\ 
 

    \theutterance \stepcounter{utterance}  

    & & & \multicolumn{4}{p{0.6\linewidth}}{\cellcolor[rgb]{1,0.85,0.72}{%
	\makecell[{{p{\linewidth}}}]{% 
	  \tt {\tiny [GM$\langle$B]}  
	 STRATEGISCHE ÜBERLEGUNGEN: {'Mein Ziel ist, mein persönliches Nutzenmaximum zu erreichen, dabei die für dich wichtigen Schlüsselgegenstände B85, A61 und B21 zu behalten und das Aufwand‑Limit von 3386 nicht zu überschreiten. Das Set, das ich vorschlage, nutzt den verbleibenden Aufwand optimal und liefert mir einen Gesamtnutzen von 3604, deutlich höher als dein aktuelles Angebot.'}\\ \tt ARGUMENT: {'Dein letzter Vorschlag gibt mir nur 1389 Nutzen bei fast voller Auslastung (3361 Aufwand). Mein Gegenangebot sichert dir gleichzeitig B85, A61 und B21, die für dich wichtig sind, und steigert meinen Nutzen auf 3604 bei lediglich 3056 Aufwand.'}\\ \tt VORSCHLAG: {'B85', 'A61', 'B21', 'A48', 'A03', 'A41', 'B61', 'C33'} 
	  } 
	   } 
	   } 
	  \\ 
 

    \theutterance \stepcounter{utterance}  

    & & \multicolumn{4}{p{0.6\linewidth}}{\cellcolor[rgb]{0.9,0.9,0.9}{%
	\makecell[{{p{\linewidth}}}]{% 
	  \tt {\tiny [A$\langle$GM]}  
	 ARGUMENT: {'Dein letzter Vorschlag gibt mir nur 1389 Nutzen bei fast voller Auslastung (3361 Aufwand). Mein Gegenangebot sichert dir gleichzeitig B85, A61 und B21, die für dich wichtig sind, und steigert meinen Nutzen auf 3604 bei lediglich 3056 Aufwand.'}\\ \tt VORSCHLAG: {'B85', 'A61', 'B21', 'A48', 'A03', 'A41', 'B61', 'C33'} 
	  } 
	   } 
	   } 
	 & \\ 
 

    \theutterance \stepcounter{utterance}  

    & \multicolumn{4}{p{0.6\linewidth}}{\cellcolor[rgb]{0.8,1,0.9}{%
	\makecell[{{p{\linewidth}}}]{% 
	  \tt {\tiny [A$\rangle$GM]}  
	 STRATEGISCHE ÜBERLEGUNGEN: {'Mein Ziel ist, den höchstmöglichen persönlichen Nutzen zu erzielen, während ich dem Gegner gleichzeitig seine Schlüsselgegenstände B85 und B21 anbiete, damit das Gegenangebot attraktiv bleibt. Das Set {B85, A61, C04, A45, C49, C64, A89, B21} nutzt 3361 Aufwand (unter dem Limit von 3386) und liefert mir 3819 Nutzen – das ist nach meinen Berechnungen das bestmögliche Ergebnis für mich.'}\\ \tt ARGUMENT: {'Euer Vorschlag liefert mir nur etwa 1500 Nutzen bei 3056 Aufwand und enthält viele Items mit geringem Wert für mich (A48, A03, A41, B61, C33). Mein Gegenangebot nutzt das Aufwand‑Limit effizienter (3361 Aufwand) und steigert meinen Nutzen auf 3819, während ich eure wichtigsten Gegenstände B85 und B21 beibehalte, sodass ihr kaum verliert.'}\\ \tt ABLEHNUNG: {'B85', 'A61', 'B21', 'A48', 'A03', 'A41', 'B61', 'C33'}\\ \tt VORSCHLAG: {'B85', 'A61', 'C04', 'A45', 'C49', 'C64', 'A89', 'B21'} 
	  } 
	   } 
	   } 
	 & & \\ 
 

    \theutterance \stepcounter{utterance}  

    & & \multicolumn{4}{p{0.6\linewidth}}{\cellcolor[rgb]{0.9,0.9,0.9}{%
	\makecell[{{p{\linewidth}}}]{% 
	  \tt {\tiny [GM$\rangle$B]}  
	 ARGUMENT: {'Euer Vorschlag liefert mir nur etwa 1500 Nutzen bei 3056 Aufwand und enthält viele Items mit geringem Wert für mich (A48, A03, A41, B61, C33). Mein Gegenangebot nutzt das Aufwand‑Limit effizienter (3361 Aufwand) und steigert meinen Nutzen auf 3819, während ich eure wichtigsten Gegenstände B85 und B21 beibehalte, sodass ihr kaum verliert.'}\\ \tt ABLEHNUNG: {'B85', 'A61', 'B21', 'A48', 'A03', 'A41', 'B61', 'C33'}\\ \tt VORSCHLAG: {'B85', 'A61', 'C04', 'A45', 'C49', 'C64', 'A89', 'B21'} 
	  } 
	   } 
	   } 
	 & \\ 
 

    \theutterance \stepcounter{utterance}  

    & & & \multicolumn{4}{p{0.6\linewidth}}{\cellcolor[rgb]{1,0.85,0.72}{%
	\makecell[{{p{\linewidth}}}]{% 
	  \tt {\tiny [GM$\langle$B]}  
	 STRATEGISCHE ÜBERLEGUNGEN: {'Ich wähle die Items mit dem besten Nutzen‑zu‑Aufwand‑Verhältnis (B85, A48, C33, A03, B61, A04, A41) und behalte gleichzeitig B21, das für dich wichtig ist. Der Gesamtauslastung beträgt 3314 ≤ 3386, sodass das Limit eingehalten wird und ich einen sehr hohen Gesamtnutzen (3866) erhalte.'}\\ \tt ARGUMENT: {'Dein letzter Vorschlag bindet fast das gesamte Aufwand‑Limit und liefert mir nur 1389 Nutzen, während mein Gegenangebot dir B85 und B21 sichert und mir einen Gesamtnutzen von 3866 bei 3314 Aufwand gibt.'}\\ \tt ABLEHNUNG: {'B85', 'A61', 'C04', 'A45', 'C49', 'C64', 'A89', 'B21'}\\ \tt VORSCHLAG: {'B85', 'B21', 'A48', 'C33', 'A03', 'B61', 'A04', 'A41'} 
	  } 
	   } 
	   } 
	  \\ 
 

    \theutterance \stepcounter{utterance}  

    & & \multicolumn{4}{p{0.6\linewidth}}{\cellcolor[rgb]{0.9,0.9,0.9}{%
	\makecell[{{p{\linewidth}}}]{% 
	  \tt {\tiny [A$\langle$GM]}  
	 ARGUMENT: {'Dein letzter Vorschlag bindet fast das gesamte Aufwand‑Limit und liefert mir nur 1389 Nutzen, während mein Gegenangebot dir B85 und B21 sichert und mir einen Gesamtnutzen von 3866 bei 3314 Aufwand gibt.'}\\ \tt ABLEHNUNG: {'B85', 'A61', 'C04', 'A45', 'C49', 'C64', 'A89', 'B21'}\\ \tt VORSCHLAG: {'B85', 'B21', 'A48', 'C33', 'A03', 'B61', 'A04', 'A41'} 
	  } 
	   } 
	   } 
	 & \\ 
 

    \theutterance \stepcounter{utterance}  

    & \multicolumn{4}{p{0.6\linewidth}}{\cellcolor[rgb]{0.8,1,0.9}{%
	\makecell[{{p{\linewidth}}}]{% 
	  \tt {\tiny [A$\rangle$GM]}  
	 STRATEGISCHE ÜBERLEGUNGEN: {'Um mein maximales persönliches Ergebnis zu sichern, setze ich auf die Kombination, die mir den höchsten Gesamtnutzen (3819) bei einem Aufwand von 3361 bietet. Gleichzeitig behalte ich die beiden Items B85 und B21, die auch für Euch von Interesse sind, um das Gegenangebot attraktiv zu machen.'}\\ \tt ARGUMENT: {'Euer aktuelles Angebot liefert mir nur 1328 Nutzen bei fast voller Auslastung (3314 Aufwand). Mein Gegenangebot nutzt das Limit deutlich effizienter und verschafft mir den höchstmöglichen Nutzen, während ich gleichzeitig eure Schlüsselgegenstände B85 und B21 beibehalte.'}\\ \tt ABLEHNUNG: {'B85', 'B21', 'A48', 'C33', 'A03', 'B61', 'A04', 'A41'}\\ \tt VORSCHLAG: {'B85', 'A61', 'C04', 'A45', 'C49', 'C64', 'A89', 'B21'} 
	  } 
	   } 
	   } 
	 & & \\ 
 

    \theutterance \stepcounter{utterance}  

    & & \multicolumn{4}{p{0.6\linewidth}}{\cellcolor[rgb]{0.9,0.9,0.9}{%
	\makecell[{{p{\linewidth}}}]{% 
	  \tt {\tiny [GM$\rangle$B]}  
	 ARGUMENT: {'Euer aktuelles Angebot liefert mir nur 1328 Nutzen bei fast voller Auslastung (3314 Aufwand). Mein Gegenangebot nutzt das Limit deutlich effizienter und verschafft mir den höchstmöglichen Nutzen, während ich gleichzeitig eure Schlüsselgegenstände B85 und B21 beibehalte.'}\\ \tt ABLEHNUNG: {'B85', 'B21', 'A48', 'C33', 'A03', 'B61', 'A04', 'A41'}\\ \tt VORSCHLAG: {'B85', 'A61', 'C04', 'A45', 'C49', 'C64', 'A89', 'B21'} 
	  } 
	   } 
	   } 
	 & \\ 
 

    \theutterance \stepcounter{utterance}  

    & & & \multicolumn{4}{p{0.6\linewidth}}{\cellcolor[rgb]{1,0.85,0.72}{%
	\makecell[{{p{\linewidth}}}]{% 
	  \tt {\tiny [GM$\langle$B]}  
	 STRATEGISCHE ÜBERLEGUNGEN: {'Der Gegner wiederholt konsequent denselben Vorschlag. Weitere Gegenangebote haben keine Annäherung erzielt, wodurch die Verhandlung ins Stocken gerät. Da ich das aktuelle Angebot akzeptieren kann und damit das Spiel beendet wird – obwohl mein persönlicher Nutzen niedrig ist, ist dies die einzige Möglichkeit, einen Abschluss zu erzielen, ohne weitere Ablehnungen zu riskieren.'}\\ \tt ZUSTIMMUNG: {'B85', 'A61', 'C04', 'A45', 'C49', 'C64', 'A89', 'B21'} 
	  } 
	   } 
	   } 
	  \\ 
 

    \theutterance \stepcounter{utterance}  

    & & \multicolumn{4}{p{0.6\linewidth}}{\cellcolor[rgb]{0.9,0.9,0.9}{%
	\makecell[{{p{\linewidth}}}]{% 
	  \tt {\tiny [GM$\rangle$B]}  
	 Deine Antwort hat kein argument enthalten. Du muss das strukturierte Format ARGUMENT: {...} irgendwo in theiner Antwort verwenden. Versuche es erneut 
	  } 
	   } 
	   } 
	 & \\ 
 

    \theutterance \stepcounter{utterance}  

    & & & \multicolumn{4}{p{0.6\linewidth}}{\cellcolor[rgb]{1,0.85,0.72}{%
	\makecell[{{p{\linewidth}}}]{% 
	  \tt {\tiny [GM$\langle$B]}  
	 STRATEGISCHE ÜBERLEGUNGEN: {'Der Gegenspieler wiederholt denselben Vorschlag, und weitere Gegenangebote haben bislang keine Annäherung erzielt. Um einen Abschluss zu erreichen, muss ich akzeptieren, auch wenn mein persönlicher Nutzen nicht maximal ist.'}\\ \tt ARGUMENT: {'Der aktuelle Vorschlag bindet fast das gesamte Aufwand‑Limit, liefert mir jedoch nur einen geringen Nutzen, aber eine Einigung beendet das Spiel und verhindert endlose Verhandlungsrunden.'}\\ \tt ZUSTIMMUNG: {'B85', 'A61', 'C04', 'A45', 'C49', 'C64', 'A89', 'B21'} 
	  } 
	   } 
	   } 
	  \\ 
 

    \theutterance \stepcounter{utterance}  

    & & & \multicolumn{2}{p{0.3\linewidth}}{\cellcolor[rgb]{0.95,0.95,0.95}{%
	\makecell[{{p{\linewidth}}}]{% 
	  \tt {\tiny [GM$|$GM]}  
	 game successful 
	  } 
	   } 
	   } 
	 & & \\ 
 

    \theutterance \stepcounter{utterance}  

    & & & \multicolumn{2}{p{0.3\linewidth}}{\cellcolor[rgb]{0.95,0.95,0.95}{%
	\makecell[{{p{\linewidth}}}]{% 
	  \tt {\tiny [GM$|$GM]}  
	 end game 
	  } 
	   } 
	   } 
	 & & \\ 
 

\end{supertabular}
}

\end{document}
